% ==========================================================================
% sotOS — Manual Técnico del Micronúcleo
% Primera Mitad: Fundamentos (Capítulos 1–4)
% ==========================================================================
\documentclass[11pt,twoside,openright]{book}

% --- Codificación y tipografía ---
\usepackage[T1]{fontenc}
\usepackage[utf8]{inputenc}
\usepackage[spanish]{babel}
\usepackage{lmodern}
\usepackage{microtype}

% --- Matemáticas ---
\usepackage{amsmath,amssymb,amsthm,mathtools}

% --- Algoritmos ---
\usepackage[ruled,vlined,linesnumbered,spanish]{algorithm2e}

% --- Gráficos ---
\usepackage{tikz}
\usetikzlibrary{arrows.meta,positioning,calc,shapes.geometric,automata,
                 decorations.pathreplacing,fit,backgrounds,chains}

% --- Código fuente ---
\usepackage{listings}
\lstset{
  language={[x86masm]Assembler},
  basicstyle=\ttfamily\small,
  keywordstyle=\bfseries\color{blue!70!black},
  commentstyle=\itshape\color{green!50!black},
  stringstyle=\color{red!60!black},
  numbers=left, numberstyle=\tiny\color{gray},
  frame=single, framesep=3pt,
  breaklines=true,
  showstringspaces=false,
  tabsize=4,
  captionpos=b,
}
\lstdefinelanguage{Rust}{
  morekeywords={fn,let,mut,pub,struct,enum,impl,self,Self,const,static,
                unsafe,extern,use,mod,match,if,else,loop,while,for,in,
                return,break,continue,where,trait,type,as,ref,true,false,
                Some,None,Option,Result,Ok,Err,u8,u16,u32,u64,usize,
                i32,i64,bool,crate},
  sensitive=true,
  morecomment=[l]{//},
  morecomment=[s]{/*}{*/},
  morestring=[b]",
}

% --- Tablas ---
\usepackage{booktabs,array,longtable}

% --- Hipervínculos ---
\usepackage[colorlinks=true,linkcolor=blue!60!black,
            citecolor=green!50!black,urlcolor=red!50!black]{hyperref}
\usepackage{cleveref}

% --- Geometría ---
\usepackage[margin=2.5cm]{geometry}

% --- Encabezados ---
\usepackage{fancyhdr}
\pagestyle{fancy}
\fancyhf{}
\fancyhead[LE]{\slshape\nouppercase{\leftmark}}
\fancyhead[RO]{\slshape\nouppercase{\rightmark}}
\fancyfoot[C]{\thepage}
\renewcommand{\headrulewidth}{0.4pt}

% --- Entornos de teoremas ---
\theoremstyle{definition}
\newtheorem{definition}{Definición}[chapter]
\newtheorem{example}{Ejemplo}[chapter]

\theoremstyle{plain}
\newtheorem{theorem}{Teorema}[chapter]
\newtheorem{lemma}[theorem]{Lema}
\newtheorem{corollary}[theorem]{Corolario}
\newtheorem{proposition}[theorem]{Proposición}

\theoremstyle{remark}
\newtheorem{remark}{Observación}[chapter]
\newtheorem{invariant}{Invariante}[chapter]

% --- Macros del proyecto ---
\newcommand{\sotos}{\textsc{sotOS}}
\newcommand{\capid}{\texttt{CapId}}
\newcommand{\poolhandle}{\texttt{PoolHandle}}
\newcommand{\bigO}{\mathcal{O}}
\newcommand{\rights}{\mathcal{R}}
\newcommand{\rightsall}{\rights_{\mathrm{ALL}}}
\newcommand{\bitand}{\mathbin{\&}}
\newcommand{\bitor}{\mathbin{|}}
\newcommand{\genmax}{G_{\max}}
\newcommand{\slotmax}{S_{\max}}
\newcommand{\imark}[1]{\textbf{#1}}

% ==========================================================================
\begin{document}

% --- Portada ---
\begin{titlepage}
\centering
\vspace*{3cm}
{\Huge\bfseries Manual Técnico de \sotos\\[0.3cm]}
{\Large Un Micronúcleo Verificable con\\Bypass Selectivo de Exonúcleo}\\[2cm]
{\large Versión 1.0 — Marzo 2026}\\[3cm]
{\Large Lucas Sotomayor}\\[1cm]
\vfill
{\small
Objetivo de diseño: $\leq 10\,000$ líneas de código en el núcleo\\
Cinco primitivas: planificación, virtualización de IRQ, IPC, capacidades, asignación de marcos
}
\end{titlepage}

% --- Tabla de contenidos ---
\frontmatter
\tableofcontents

% ====================================================================
\mainmatter
\part{Fundamentos}
% ====================================================================

% ====================================================================
\chapter{Introducción y Filosofía de Diseño}
\label{chap:intro}
% ====================================================================

\section{Motivación}

La historia de los sistemas operativos está marcada por una tensión fundamental entre
\emph{rendimiento} y \emph{seguridad}. Los núcleos monolíticos como Linux incorporan
millones de líneas de código en el espacio privilegiado (Ring~0), incluyendo controladores
de dispositivos, sistemas de archivos y pilas de red. Cada línea de código en Ring~0
es una línea que, ante un error, puede comprometer la integridad completa del sistema.

\sotos{} adopta una posición radical: \emph{el código en Ring~0 debe ser mínimo y
verificable}. Con un objetivo de $\leq 10\,000$ líneas de código (LOC) en el núcleo,
el sistema implementa exactamente cinco primitivas fundamentales:

\begin{enumerate}
    \item \textbf{Planificación}: Colas multinivel (MLQ) con robo de trabajo entre CPUs.
    \item \textbf{Virtualización de IRQ}: Delegación de interrupciones a espacio de usuario.
    \item \textbf{IPC}: Endpoints síncronos (estilo L4) y canales asíncronos (ring buffers).
    \item \textbf{Capacidades}: Seguridad basada en tokens opacos con derechos asociados.
    \item \textbf{Asignación de marcos}: Bitmap de marcos físicos de 4\,KiB.
\end{enumerate}

Todo lo demás --- memoria virtual (paginación), sistema de archivos, controladores,
pila de red --- se ejecuta en espacio de usuario como servidores independientes.

\section{Principios de Diseño}

\subsection{Base Computacional de Confianza Mínima (TCB)}

\begin{definition}[Base Computacional de Confianza]
\label{def:tcb}
La \emph{Base Computacional de Confianza} (TCB, por \emph{Trusted Computing Base}) de un
sistema es el conjunto mínimo de componentes de hardware, firmware y software cuyo
funcionamiento correcto es suficiente para garantizar la política de seguridad del sistema.
\end{definition}

En un núcleo monolítico, la TCB incluye todo el código del kernel, típicamente millones
de LOC. En \sotos{}, la TCB comprende únicamente:

\begin{itemize}
    \item El micronúcleo ($\leq 10\,000$ LOC).
    \item El hardware (CPU x86\_64, controlador de memoria).
    \item El cargador de arranque (Limine v8).
\end{itemize}

\begin{theorem}[Reducción de superficie de ataque]
\label{thm:tcb-reduction}
Sea $n_M$ el número de líneas de código del kernel monolítico y $n_\mu$ el del micronúcleo,
con $n_\mu \ll n_M$. Si la tasa de defectos es $\lambda$ defectos por línea (constante),
entonces el número esperado de defectos en la TCB se reduce de $\lambda n_M$ a
$\lambda n_\mu$, logrando una reducción proporcional a $n_\mu / n_M$.
\end{theorem}

\begin{proof}
Sea $D(n) = \lambda n$ el número esperado de defectos en $n$ líneas de código, donde
$\lambda > 0$ es la densidad de defectos. Entonces:
\[
    \frac{D(n_\mu)}{D(n_M)} = \frac{\lambda n_\mu}{\lambda n_M} = \frac{n_\mu}{n_M}.
\]
Para \sotos{} con $n_\mu \leq 10\,000$ y un kernel Linux típico con $n_M \approx 30\,000\,000$:
\[
    \frac{n_\mu}{n_M} \leq \frac{10\,000}{30\,000\,000} \approx 3.3 \times 10^{-4}.
\]
Esto implica una reducción de más de tres órdenes de magnitud en el número esperado de
defectos dentro de la TCB.
\end{proof}

\subsection{Seguridad Basada en Capacidades}

A diferencia de los modelos de control de acceso discrecional (DAC) o mandatorio (MAC)
que asocian permisos a identidades o etiquetas, \sotos{} emplea un modelo de
\emph{capacidades} (\emph{capabilities}): tokens opacos que simultáneamente identifican
un objeto del kernel y codifican los derechos de acceso sobre él.

\begin{definition}[Capacidad]
\label{def:capability}
Una \emph{capacidad} es un par $c = (o, r)$ donde $o \in \mathcal{O}$ es un objeto del
kernel y $r \in 2^{\rights}$ es un conjunto de derechos, con
$\rights = \{\textsc{Read}, \textsc{Write}, \textsc{Execute}, \textsc{Grant}, \textsc{Revoke}\}$.
\end{definition}

Este modelo elimina las vulnerabilidades de tipo \emph{confused deputy}, donde un servicio
privilegiado es engañado para abusar de sus permisos, porque el derecho de acceso viaja
con la referencia al objeto y no con la identidad del invocador.

\subsection{Bypass Selectivo de Exonúcleo}

Un exonúcleo expone directamente los recursos de hardware al espacio de usuario,
eliminando abstracciones del kernel. \sotos{} adopta un enfoque híbrido: el micronúcleo
provee abstracciones seguras (IPC, capacidades, planificación), pero ciertos caminos
críticos de rendimiento permiten acceso directo a hardware mediado por capacidades de
tipo \texttt{IoPort} y \texttt{Memory}.

\begin{example}
Un controlador de disco NVMe en espacio de usuario recibe una capacidad
$c_{\text{mmio}} = (\text{Memory}(\texttt{0xFE000000}, 4096), \{\textsc{Read}, \textsc{Write}\})$
que le permite acceder directamente a los registros MMIO del dispositivo sin pasar por
el kernel en cada operación de E/S.
\end{example}

\section{Comparación Arquitectónica}

\begin{table}[ht]
\centering
\caption{Comparación de arquitecturas de sistemas operativos.}
\label{tab:comparison}
\begin{tabular}{@{}lccc@{}}
\toprule
\textbf{Propiedad} & \textbf{Monolítico} & \textbf{Micronúcleo} & \textbf{Exonúcleo} \\
\midrule
Código en Ring~0        & Todo         & Mínimo        & Mínimo \\
TCB (LOC)               & $\sim 10^7$  & $\sim 10^4$   & $\sim 10^3$ \\
Abstracción de HW       & Completa     & Parcial       & Ninguna \\
Rendimiento de IPC      & N/A (in-kernel) & Crítico    & N/A \\
Aislamiento de fallos   & Ninguno      & Por proceso   & Por proceso \\
Controladores           & En kernel    & En usuario    & En usuario \\
Verificabilidad         & Impráctica   & Factible      & Factible \\
\bottomrule
\end{tabular}
\end{table}

\sotos{} se ubica entre el micronúcleo puro y el exonúcleo: provee las cinco primitivas
del micronúcleo con la opción de bypass directo a hardware cuando la política de seguridad
lo permite (vía capacidades \texttt{IoPort} y \texttt{Memory}).

\section{Arquitectura General del Sistema}

\begin{figure}[ht]
\centering
\begin{tikzpicture}[
    >=Stealth,
    box/.style={draw, rounded corners=3pt, minimum width=2.8cm, minimum height=0.9cm,
                font=\small, align=center},
    kernbox/.style={box, fill=red!15, draw=red!50!black},
    userbox/.style={box, fill=blue!12, draw=blue!50!black},
    hwbox/.style={box, fill=gray!20, draw=gray!60!black},
    ipc/.style={->, thick, dashed, blue!60!black},
    syscall/.style={->, thick, red!60!black},
    hw/.style={<->, thick, gray!50!black},
]
% Hardware layer
\node[hwbox] (cpu) at (0, 0) {CPU\\x86\_64};
\node[hwbox] (mem) at (3.5, 0) {DRAM\\DDR4};
\node[hwbox] (disk) at (7, 0) {NVMe\\SSD};
\node[hwbox] (nic) at (10.5, 0) {NIC\\virtio};

% Kernel layer
\node[kernbox] (sched) at (0, 2.2) {Planificador\\MLQ+WS};
\node[kernbox] (ipc) at (3.5, 2.2) {IPC\\Endpoints};
\node[kernbox] (caps) at (7, 2.2) {Capacidades\\CDT};
\node[kernbox] (frame) at (10.5, 2.2) {Asignador\\de Marcos};

% Ring boundary
\draw[very thick, red!40!black, dashed] (-2, 1.4) -- (12.5, 1.4)
    node[right, font=\footnotesize\itshape] {Ring 0 / Ring 3};

% User layer
\node[userbox] (vmm) at (0, 4) {VMM\\(Paginación)};
\node[userbox] (fs) at (3.5, 4) {FS\\(VFS+ObjStore)};
\node[userbox] (nvme) at (7, 4) {Driver\\NVMe};
\node[userbox] (net) at (10.5, 4) {Pila de\\Red};

% Shell
\node[userbox, fill=green!12, draw=green!50!black] (shell) at (5.25, 5.8)
    {Shell LUCAS\\(Intérprete)};

% Arrows: syscall
\draw[syscall] (vmm) -- (sched);
\draw[syscall] (fs) -- (ipc);
\draw[syscall] (nvme) -- (caps);
\draw[syscall] (net) -- (frame);

% Arrows: IPC
\draw[ipc] (shell) -- (fs);
\draw[ipc] (shell) -- (net);
\draw[ipc] (fs.south east) -- (nvme.north west);
\draw[ipc] (vmm.east) -- (fs.west);

% HW arrows
\draw[hw] (cpu) -- (sched);
\draw[hw] (mem) -- (frame);
\draw[hw] (disk) -- (nvme);
\draw[hw] (nic) -- (net);

% Legend
\matrix[draw, column sep=8pt, row sep=2pt, anchor=south west,
        font=\footnotesize] at (-2, -1.2) {
    \node[kernbox, minimum width=1cm, minimum height=0.4cm] {}; &
    \node {Kernel (Ring 0)}; &
    \node[userbox, minimum width=1cm, minimum height=0.4cm] {}; &
    \node {Usuario (Ring 3)}; \\
};

\end{tikzpicture}
\caption{Arquitectura de capas de \sotos{}: hardware, micronúcleo (Ring~0) y
servidores de espacio de usuario (Ring~3). Las flechas continuas representan
llamadas al sistema; las discontinuas representan IPC entre servidores.}
\label{fig:architecture}
\end{figure}

\section{Estructura del Código Fuente}

El repositorio sigue una organización modular:

\begin{description}
    \item[\texttt{kernel/}] El micronúcleo (Ring~0):
    \begin{description}
        \item[\texttt{arch/x86\_64/}] Código específico de CPU: serial, GDT, IDT, LAPIC.
        \item[\texttt{mm/}] Asignador de marcos físicos (bitmap).
        \item[\texttt{cap/}] Tabla de capacidades y CDT.
        \item[\texttt{ipc/}] Endpoints (síncronos) y canales (asíncronos).
        \item[\texttt{sched/}] Gestión de hilos y planificador.
        \item[\texttt{pool.rs}] Asignador genérico \texttt{Pool<T>} con contadores de generación.
    \end{description}
    \item[\texttt{libs/sotos-common/}] ABI compartida kernel/usuario (syscalls, errores).
    \item[\texttt{services/}] Servidores de usuario: VMM, FS, controladores, shell.
    \item[\texttt{boot/}] Configuración del cargador Limine.
    \item[\texttt{scripts/}] Herramientas de construcción (\texttt{mkimage.py}).
\end{description}

\section{Notación y Convenciones}

A lo largo de este manual empleamos la siguiente notación:

\begin{itemize}
    \item $\mathcal{O}$ denota el conjunto de objetos del kernel.
    \item $\rights = \{\textsc{R}, \textsc{W}, \textsc{X}, \textsc{G}, \textsc{Rv}\}$ es el
          conjunto universal de derechos (Read, Write, Execute, Grant, Revoke).
    \item Para un entero $x$ de $n$ bits, $x[a{:}b]$ denota los bits desde la posición $b$
          (inclusive) hasta la posición $a$ (inclusive), con el bit 0 siendo el menos significativo.
    \item $\lfloor x \rfloor_n$ denota $x \bmod n$.
    \item $|S|$ denota la cardinalidad de un conjunto finito $S$.
    \item Los fragmentos de código se presentan en Rust (el lenguaje de implementación).
\end{itemize}


% ====================================================================
\chapter{Sistema de Capacidades}
\label{chap:capabilities}
% ====================================================================

El sistema de capacidades es el pilar de seguridad de \sotos{}. Cada objeto del kernel ---
marcos de memoria, endpoints de IPC, hilos, puertos de E/S --- solo es accesible a través
de una capacidad que simultáneamente identifica el objeto y codifica los derechos de acceso.
No existe un mecanismo alternativo: ni identificadores de proceso, ni listas de control de
acceso, ni anillos de protección adicionales. Las capacidades son \emph{el único} mecanismo
de autorización.

\section{Definiciones Formales}

\begin{definition}[Universo de Objetos]
\label{def:object-universe}
El \emph{universo de objetos} $\mathcal{O}$ de \sotos{} es la unión disjunta:
\[
    \mathcal{O} = \mathcal{O}_{\text{Mem}} \cup \mathcal{O}_{\text{Ep}} \cup
    \mathcal{O}_{\text{Ch}} \cup \mathcal{O}_{\text{Thr}} \cup
    \mathcal{O}_{\text{Irq}} \cup \mathcal{O}_{\text{Io}} \cup
    \mathcal{O}_{\text{Notif}} \cup \mathcal{O}_{\text{Dom}} \cup
    \mathcal{O}_{\text{AS}}
\]
donde cada subconjunto corresponde a un tipo de objeto:
\begin{itemize}
    \item $\mathcal{O}_{\text{Mem}} = \{(\text{base}, \text{size}) \mid \text{base}, \text{size} \in \mathbb{N}_{64}\}$:
          regiones de memoria física.
    \item $\mathcal{O}_{\text{Ep}} = \{e_i \mid i \in [0, C \cdot E)\}$: endpoints de IPC,
          donde $C = \texttt{MAX\_CPUS}$ y $E = \texttt{MAX\_EP\_PER\_CORE} = 64$.
    \item $\mathcal{O}_{\text{Ch}} = \{ch_i\}$: canales asíncronos.
    \item $\mathcal{O}_{\text{Thr}} = \{t_i\}$: hilos.
    \item $\mathcal{O}_{\text{Irq}} = \{\text{irq}_n \mid n \in [0, 256)\}$: líneas de interrupción.
    \item $\mathcal{O}_{\text{Io}} = \{(\text{base}, \text{count}) \mid \text{base} \in [0, 2^{16}), \text{count} \in [1, 2^{16}]\}$:
          rangos de puertos x86 de E/S.
    \item $\mathcal{O}_{\text{Notif}} = \{n_i\}$: objetos de notificación (semáforos binarios).
    \item $\mathcal{O}_{\text{Dom}} = \{d_i\}$: dominios de planificación.
    \item $\mathcal{O}_{\text{AS}} = \{(\text{cr3})\}$: espacios de direcciones (PML4).
\end{itemize}
\end{definition}

\begin{definition}[Derechos]
\label{def:rights}
El conjunto de derechos es $\rights = \{\textsc{R}, \textsc{W}, \textsc{X}, \textsc{G}, \textsc{Rv}\}$,
codificado como un bitmask de 5 bits:
\[
    \textsc{R} = 2^0, \quad
    \textsc{W} = 2^1, \quad
    \textsc{X} = 2^2, \quad
    \textsc{G} = 2^3, \quad
    \textsc{Rv} = 2^4.
\]
El derecho universal es $\rightsall = \textsc{R} \bitor \textsc{W} \bitor \textsc{X} \bitor \textsc{G} \bitor \textsc{Rv} = \texttt{0x1F}$.
Un subconjunto de derechos $r \subseteq \rights$ se representa como un entero $r \in [0, 31]$.
\end{definition}

\begin{definition}[Capacidad]
\label{def:cap-formal}
Una \emph{capacidad} es una terna $c = (o, r, p)$ donde:
\begin{itemize}
    \item $o \in \mathcal{O}$ es el objeto referenciado.
    \item $r \in 2^{\rights}$ es el conjunto de derechos (bitmask de 5 bits).
    \item $p \in \{\bot\} \cup \mathcal{C}$ es la capacidad padre en el árbol de derivación
          ($\bot$ indica capacidad raíz).
\end{itemize}
donde $\mathcal{C}$ es el conjunto de todas las capacidades existentes en el sistema.
\end{definition}

\section{Pool Genérico con Contadores de Generación}
\label{sec:pool}

Antes de describir la tabla de capacidades, definimos la estructura de datos subyacente
que proporciona asignación $\bigO(1)$ con protección contra el problema ABA.

\subsection{El Problema ABA}

\begin{definition}[Problema ABA]
\label{def:aba}
En un sistema con reciclaje de identificadores, el \emph{problema ABA} ocurre cuando:
\begin{enumerate}
    \item Un actor obtiene una referencia al objeto $A$ en el slot $i$.
    \item El objeto $A$ es liberado del slot $i$.
    \item Un nuevo objeto $B$ es asignado al slot $i$.
    \item El actor usa su referencia obsoleta, que ahora apunta a $B$ en lugar de $A$.
\end{enumerate}
\end{definition}

En el contexto de un kernel, este problema es particularmente peligroso: si un proceso
mantiene un \capid{} cuyo slot fue reciclado, podría inadvertidamente acceder a un
objeto completamente diferente con derechos que no le corresponden.

\subsection{Codificación de Handles}

\sotos{} resuelve el problema ABA mediante \emph{contadores de generación} empaquetados
en cada handle.

\begin{definition}[PoolHandle]
\label{def:poolhandle}
Un \poolhandle{} es un entero sin signo de 32 bits con la siguiente codificación:
\[
    h = (g \ll 20) \bitor s
\]
donde:
\begin{itemize}
    \item $s = h[19{:}0]$ es el \emph{índice de slot}, con $s \in [0, 2^{20})$
          (máximo 1\,048\,575 slots).
    \item $g = h[31{:}20]$ es el \emph{contador de generación}, con $g \in [0, 2^{12})$
          (wrap-around a 4096).
\end{itemize}
\end{definition}

\begin{figure}[ht]
\centering
\begin{tikzpicture}[
    bit/.style={draw, minimum width=0.32cm, minimum height=0.6cm, font=\tiny},
    label/.style={font=\small\bfseries},
]
% Generation bits (31:20)
\foreach \i in {0,...,11} {
    \node[bit, fill=orange!25] (g\i) at (\i*0.32, 0) {};
}
% Index bits (19:0)
\foreach \i in {0,...,19} {
    \node[bit, fill=blue!15] (s\i) at ({(12+\i)*0.32}, 0) {};
}

% Labels
\node[label, above=3pt of g0, anchor=south west] {31};
\node[label, above=3pt of g11, anchor=south east] {20};
\node[label, above=3pt of s0, anchor=south west] {19};
\node[label, above=3pt of s19, anchor=south east] {0};

% Braces
\draw[decorate, decoration={brace, amplitude=5pt, mirror}]
    ([yshift=-2pt]g0.south west) -- ([yshift=-2pt]g11.south east)
    node[midway, below=6pt, font=\small] {generación (12 bits)};
\draw[decorate, decoration={brace, amplitude=5pt, mirror}]
    ([yshift=-2pt]s0.south west) -- ([yshift=-2pt]s19.south east)
    node[midway, below=6pt, font=\small] {índice de slot (20 bits)};
\end{tikzpicture}
\caption{Codificación de un \poolhandle{}: los 12 bits superiores almacenan el
contador de generación y los 20 bits inferiores el índice del slot.}
\label{fig:poolhandle}
\end{figure}

\subsection{Operaciones del Pool}

El \texttt{Pool<T>} soporta tres operaciones fundamentales:

\begin{algorithm}[H]
\caption{\textsc{Pool-Alloc}$(P, v)$: Asignar un valor $v$ en el pool $P$.}
\label{alg:pool-alloc}
\KwIn{Pool $P$, valor $v$ de tipo $T$}
\KwOut{Handle $h$ con generación correcta}
\If{$P.\text{free\_list} \neq \emptyset$}{
    $i \gets P.\text{free\_list}.\textsc{Pop}()$\;
    $P.\text{slots}[i] \gets \textsc{Some}(v)$\;
    \Return $\textsc{Pack}(i, P.\text{gens}[i])$\;
}
$i \gets |P.\text{slots}|$\;
$P.\text{slots}.\textsc{Push}(\textsc{Some}(v))$\;
$P.\text{gens}.\textsc{Push}(0)$\;
\Return $\textsc{Pack}(i, 0)$\;
\end{algorithm}

\begin{algorithm}[H]
\caption{\textsc{Pool-Free}$(P, h)$: Liberar el slot referenciado por $h$.}
\label{alg:pool-free}
\KwIn{Pool $P$, handle $h$}
\KwOut{$\textsc{Some}(v)$ si el handle era válido, $\textsc{None}$ en caso contrario}
$(i, g) \gets \textsc{Unpack}(h)$\;
\If{$i \geq |P.\text{slots}|$ \textbf{or} $P.\text{gens}[i] \neq g$}{
    \Return $\textsc{None}$\tcp*{Handle obsoleto o fuera de rango}
}
\If{$P.\text{slots}[i] = \textsc{None}$}{
    \Return $\textsc{None}$\tcp*{Doble liberación}
}
$v \gets P.\text{slots}[i].\textsc{Take}()$\;
$P.\text{gens}[i] \gets (P.\text{gens}[i] + 1) \bmod 2^{12}$\;
$P.\text{free\_list}.\textsc{Push}(i)$\;
\Return $\textsc{Some}(v)$\;
\end{algorithm}

\begin{algorithm}[H]
\caption{\textsc{Pool-Get}$(P, h)$: Buscar el valor referenciado por $h$.}
\label{alg:pool-get}
\KwIn{Pool $P$, handle $h$}
\KwOut{$\textsc{Some}(\&v)$ si válido, $\textsc{None}$ en caso contrario}
$(i, g) \gets \textsc{Unpack}(h)$\;
\If{$i \geq |P.\text{slots}|$ \textbf{or} $P.\text{gens}[i] \neq g$}{
    \Return $\textsc{None}$\;
}
\Return $P.\text{slots}[i].\textsc{AsRef}()$\;
\end{algorithm}

\subsection{Corrección de la Protección ABA}

\begin{theorem}[El contador de generación previene uso-después-de-liberar]
\label{thm:aba}
Sea $P$ un \texttt{Pool<T>} y sea $h = \textsc{Pack}(i, g)$ un handle obtenido de
\textsc{Pool-Alloc}. Si el slot $i$ es liberado y reasignado $k$ veces antes de que
$h$ sea usado en \textsc{Pool-Get}, entonces $h$ será rechazado siempre que
$k \not\equiv 0 \pmod{2^{12}}$.
\end{theorem}

\begin{proof}
Cada invocación de \textsc{Pool-Free} incrementa $P.\text{gens}[i]$ módulo $2^{12}$.
Después de $k$ liberaciones, la generación almacenada es:
\[
    g' = (g + k) \bmod 2^{12}.
\]
La verificación en \textsc{Pool-Get} compara $g$ (la generación del handle) con $g'$
(la generación actual del slot). La comparación falla si y solo si $g \neq g'$, es decir,
$k \not\equiv 0 \pmod{2^{12}}$.

Para $\genmax = 2^{12} = 4096$, un handle obsoleto será correctamente rechazado en las
primeras 4095 reasignaciones del mismo slot. Solo si el slot es reciclado exactamente
un múltiplo de 4096 veces entre la obtención y el uso del handle, este sería incorrectamente
aceptado --- una condición estadísticamente negligible en la práctica.
\end{proof}

\begin{corollary}[Probabilidad de colisión ABA]
\label{cor:aba-prob}
Si los accesos a handles obsoletos son uniformemente aleatorios con respecto al número de
reciclajes del slot, la probabilidad de que un handle obsoleto sea aceptado es
$1/\genmax = 1/4096 \approx 2.44 \times 10^{-4}$ por intento.
\end{corollary}

\begin{remark}
En la práctica, la ventana de vulnerabilidad es aún menor porque:
(a) los slots se reciclan mediante una lista libre LIFO, distribuyendo la reutilización
entre todos los slots disponibles; y
(b) un proceso que intenta usar un handle obsoleto típicamente lo hace después de pocas
reasignaciones, no miles.
\end{remark}

\section{Tabla de Capacidades}

La tabla de capacidades es la estructura central de seguridad del kernel. Almacena todas
las capacidades del sistema, protegida por un spinlock global.

\begin{definition}[Tabla de Capacidades]
\label{def:captable}
La \emph{tabla de capacidades} $\mathcal{T}$ es un \texttt{Pool<CapEntry>} donde cada
entrada es una terna $(o, r, p)$ según la \cref{def:cap-formal}. Las operaciones son:
\begin{itemize}
    \item $\textsc{Insert}(o, r, p) \to \capid{}$: crear una nueva capacidad.
    \item $\textsc{Validate}(\capid{}, r_{\text{req}}) \to \text{Result}$: verificar existencia
          y derechos suficientes.
    \item $\textsc{Lookup}(\capid{}) \to (o, r)$: obtener el objeto y derechos.
    \item $\textsc{Revoke}(\capid{})$: revocar la capacidad y todas sus derivadas.
\end{itemize}
\end{definition}

\subsection{Validación}

La validación es el camino crítico de rendimiento: cada llamada al sistema que opera sobre
un objeto del kernel primero valida la capacidad proporcionada por el espacio de usuario.

\begin{algorithm}[H]
\caption{\textsc{Cap-Validate}$(\mathcal{T}, h, r_{\text{req}})$}
\label{alg:cap-validate}
\KwIn{Tabla $\mathcal{T}$, handle $h$ (del espacio de usuario), derechos requeridos $r_{\text{req}}$}
\KwOut{Objeto $o$ si la validación tiene éxito, error en caso contrario}
$e \gets \textsc{Pool-Get}(\mathcal{T}, h)$\;
\lIf{$e = \textsc{None}$}{\Return $\textsc{Err}(\text{InvalidCap})$}
\If{$e.r \bitand r_{\text{req}} \neq r_{\text{req}}$}{
    \Return $\textsc{Err}(\text{NoRights})$\tcp*{Derechos insuficientes}
}
\Return $\textsc{Ok}(e.o)$\;
\end{algorithm}

\begin{proposition}[Complejidad de la validación]
\label{prop:validate-complexity}
\textsc{Cap-Validate} ejecuta en $\bigO(1)$: una búsqueda indexada en el pool
(verificación de generación + acceso a array) seguida de una operación bitwise AND.
\end{proposition}

\section{Delegación de Capacidades}

La delegación (\emph{grant}) permite a un proceso transferir acceso a un objeto a otro
proceso, con derechos opcionalmente restringidos.

\begin{definition}[Operación de Delegación]
\label{def:grant}
Dada una capacidad fuente $c_s = (o, r_s, p_s)$ con $\textsc{G} \in r_s$,
y una máscara de derechos $m \in 2^{\rights}$, la \emph{delegación} produce una nueva
capacidad:
\[
    c_d = (o, r_s \bitand m, c_s)
\]
donde $c_s$ se convierte en el padre de $c_d$ en el Árbol de Derivación de Capacidades.
\end{definition}

\begin{lemma}[Monotonía de Derechos]
\label{lem:rights-monotone}
Para cualquier capacidad derivada $c_d$ de una capacidad fuente $c_s$ mediante delegación
con máscara $m$:
\[
    r_d = r_s \bitand m \subseteq r_s.
\]
Es decir, una capacidad derivada nunca puede tener más derechos que su padre.
\end{lemma}

\begin{proof}
Para cualquier bit $b \in \{0, 1, 2, 3, 4\}$:
\[
    r_d[b] = r_s[b] \bitand m[b] \leq r_s[b]
\]
ya que $x \bitand y \leq x$ para todo $x, y \in \{0, 1\}$. Por lo tanto,
$r_d \subseteq r_s$ como conjuntos de derechos.
\end{proof}

\begin{corollary}[Cadena de Delegaciones]
\label{cor:delegation-chain}
Si $c_0 \to c_1 \to \cdots \to c_k$ es una cadena de delegaciones, entonces
$r_k \subseteq r_{k-1} \subseteq \cdots \subseteq r_0$. Los derechos son
monótonamente no crecientes a lo largo de la cadena.
\end{corollary}

\begin{proof}
Aplicar el \cref{lem:rights-monotone} inductivamente: si $r_i \subseteq r_{i-1}$ para
todo $i \leq k$ (base: $r_1 = r_0 \bitand m_1 \subseteq r_0$), entonces
$r_{i+1} = r_i \bitand m_{i+1} \subseteq r_i \subseteq r_{i-1}$.
\end{proof}

\section{Árbol de Derivación de Capacidades (CDT)}
\label{sec:cdt}

El CDT registra la relación padre-hijo entre capacidades, permitiendo la revocación
transitiva: al revocar una capacidad, se revocan automáticamente todas las capacidades
derivadas de ella.

\begin{definition}[CDT]
\label{def:cdt}
El \emph{Árbol de Derivación de Capacidades} (CDT) es un bosque dirigido
$G = (\mathcal{C}, E)$ donde:
\begin{itemize}
    \item $\mathcal{C}$ es el conjunto de capacidades vivas.
    \item $(c_p, c_d) \in E$ si $c_d$ fue creada por delegación desde $c_p$.
    \item Las capacidades raíz (insertadas directamente por el kernel) tienen $p = \bot$.
\end{itemize}
\end{definition}

\begin{figure}[ht]
\centering
\begin{tikzpicture}[
    >=Stealth,
    cap/.style={draw, circle, minimum size=1.2cm, font=\small},
    root/.style={cap, fill=red!20, draw=red!50!black, thick},
    child/.style={cap, fill=blue!15, draw=blue!50!black},
    dead/.style={cap, fill=gray!30, draw=gray, dashed, text=gray},
    edge/.style={->, thick},
    redge/.style={->, thick, red!60!black, dashed},
]
% Root capability
\node[root] (c0) at (0, 0) {$c_0$\\{\tiny RWG}};

% Level 1
\node[child] (c1) at (-3, -2) {$c_1$\\{\tiny RW}};
\node[child] (c2) at (0, -2) {$c_2$\\{\tiny RG}};
\node[child] (c3) at (3, -2) {$c_3$\\{\tiny R}};

% Level 2
\node[child] (c4) at (-4, -4) {$c_4$\\{\tiny R}};
\node[child] (c5) at (-2, -4) {$c_5$\\{\tiny RW}};
\node[child] (c6) at (0, -4) {$c_6$\\{\tiny R}};

% Edges
\draw[edge] (c0) -- (c1);
\draw[edge] (c0) -- (c2);
\draw[edge] (c0) -- (c3);
\draw[edge] (c1) -- (c4);
\draw[edge] (c1) -- (c5);
\draw[edge] (c2) -- (c6);

% Revocation illustration
\node[dead] (c1r) at (-3, -2) {$c_1$\\{\tiny RW}};
\node[dead] (c4r) at (-4, -4) {$c_4$\\{\tiny R}};
\node[dead] (c5r) at (-2, -4) {$c_5$\\{\tiny RW}};

% Revocation arrow
\draw[redge, bend left=40] ($(c0.south west)+(-0.5, -0.2)$) to
    node[left, font=\footnotesize\itshape, text=red!60!black] {\textsc{Revoke}($c_1$)}
    ($(c1.north west)+(-0.3, 0.2)$);

% Legend
\node[anchor=north west, font=\footnotesize] at (4.5, 0) {
    \begin{tabular}{@{}ll@{}}
        \tikz\node[root, minimum size=0.5cm] {}; & Raíz \\
        \tikz\node[child, minimum size=0.5cm] {}; & Derivada \\
        \tikz\node[dead, minimum size=0.5cm] {}; & Revocada \\
    \end{tabular}
};

\end{tikzpicture}
\caption{Árbol de Derivación de Capacidades. Al revocar $c_1$, se revocan transitivamente
$c_4$ y $c_5$. Las capacidades $c_2$, $c_3$ y $c_6$ no se ven afectadas.}
\label{fig:cdt}
\end{figure}

\subsection{Algoritmo de Revocación}

La revocación en \sotos{} se implementa mediante un algoritmo iterativo de punto fijo que
elimina huérfanos hasta que no quedan más.

\begin{algorithm}[H]
\caption{\textsc{Revoke}$(\mathcal{T}, h)$: Revocar una capacidad y todas sus derivadas.}
\label{alg:revoke}
\KwIn{Tabla $\mathcal{T}$, handle $h$ de la capacidad a revocar}
\textsc{Pool-Free}$(\mathcal{T}, h)$\tcp*{Liberar la capacidad objetivo}
\Repeat{$\text{orphans} = \emptyset$}{
    $\text{orphans} \gets \emptyset$\;
    \ForEach{$(h_i, e_i) \in \mathcal{T}.\text{iter}()$}{
        \If{$e_i.p \neq \bot$ \textbf{and} $\textsc{Pool-Get}(\mathcal{T}, e_i.p) = \textsc{None}$}{
            $\text{orphans} \gets \text{orphans} \cup \{h_i\}$\;
        }
    }
    \ForEach{$h_i \in \text{orphans}$}{
        \textsc{Pool-Free}$(\mathcal{T}, h_i)$\;
    }
}
\end{algorithm}

\begin{theorem}[Terminación de la Revocación]
\label{thm:revoke-terminates}
El algoritmo \textsc{Revoke} termina en a lo sumo $d + 1$ iteraciones del bucle externo,
donde $d$ es la profundidad máxima del subárbol enraizado en la capacidad revocada.
\end{theorem}

\begin{proof}
Sea $T_h$ el subárbol del CDT enraizado en $h$ (incluyendo $h$). Definimos:
\begin{itemize}
    \item $L_0 = \{h\}$ (nivel 0, la raíz).
    \item $L_k = \{c \in T_h \mid \text{dist}(h, c) = k\}$ (nivel $k$).
\end{itemize}

\noindent\textbf{Iteración 0}: Se libera $h$, lo que deja a todos los hijos directos $L_1$
como huérfanos (su padre $h$ ya no existe en el pool).

\noindent\textbf{Iteración $k$ ($k \geq 1$)}: El bucle identifica como huérfanos todos los nodos
en $L_k$ (sus padres en $L_{k-1}$ fueron liberados en la iteración anterior). Se liberan
todos los nodos de $L_k$.

\noindent\textbf{Iteración $d + 1$}: Todos los nodos de $L_d$ (el nivel más profundo) fueron
liberados en la iteración $d$. No quedan huérfanos. El conjunto $\text{orphans}$ es vacío
y el bucle termina.

Como $d \leq |\mathcal{C}|$ (profundidad máxima acotada por el número total de capacidades,
que a su vez está acotado por $\slotmax = 2^{20}$), el algoritmo termina en tiempo finito.
\end{proof}

\begin{theorem}[Corrección de la Revocación]
\label{thm:revoke-correct}
Después de ejecutar \textsc{Revoke}$(\mathcal{T}, h)$, para todo $c$ que era descendiente
de $h$ en el CDT (o $h$ mismo), $\textsc{Pool-Get}(\mathcal{T}, c) = \textsc{None}$.
\end{theorem}

\begin{proof}
Por inducción sobre la profundidad $k$ del nodo en el subárbol $T_h$.

\noindent\textbf{Caso base} ($k = 0$): $h$ es liberado directamente en la primera línea del algoritmo.

\noindent\textbf{Paso inductivo}: Supongamos que todo nodo a profundidad $\leq k$ ha sido
liberado. Sea $c$ un nodo a profundidad $k + 1$. Su padre está a profundidad $k$ y, por
hipótesis inductiva, ha sido liberado. Entonces $\textsc{Pool-Get}(\mathcal{T}, c.\text{parent})
= \textsc{None}$, por lo que $c$ es identificado como huérfano en la siguiente iteración
del bucle y es liberado.

Ningún nodo fuera de $T_h$ es afectado porque: (a) los nodos raíz ($p = \bot$) nunca
son identificados como huérfanos; (b) los nodos cuyo padre pertenece a otro subárbol
mantienen su padre vivo.
\end{proof}

\begin{proposition}[Complejidad de la Revocación]
\label{prop:revoke-complexity}
Sea $n = |\mathcal{C}|$ el número total de capacidades y $m = |T_h|$ el tamaño del
subárbol revocado. La complejidad es $\bigO(d \cdot n)$ donde $d$ es la profundidad
del subárbol, o equivalentemente $\bigO(n \cdot m)$ en el peor caso ($d = m$, cadena lineal).
\end{proposition}

\begin{remark}
Una implementación alternativa con punteros explícitos padre$\to$hijos permitiría
revocación en $\bigO(m)$, pero requeriría listas enlazadas dinámicas en cada entrada.
\sotos{} prioriza la simplicidad del slot fijo para mantener el código del kernel dentro
del objetivo de $10\,000$ LOC.
\end{remark}

\section{Tipos de Objetos y su Semántica}

\begin{table}[ht]
\centering
\caption{Tipos de objetos del kernel y operaciones válidas por derecho.}
\label{tab:cap-objects}
\begin{tabular}{@{}llp{7cm}@{}}
\toprule
\textbf{Tipo} & \textbf{Variante Rust} & \textbf{Operaciones} \\
\midrule
Memoria & \texttt{Memory\{base, size\}} &
    \textsc{R}: mapear como solo-lectura;
    \textsc{W}: mapear como lectura-escritura;
    \textsc{X}: mapear como ejecutable (W$\oplus$X enforced) \\
Endpoint & \texttt{Endpoint\{id\}} &
    \textsc{R}: recv;
    \textsc{W}: send;
    \textsc{G}: transferir cap en IPC \\
Canal & \texttt{Channel\{id\}} &
    \textsc{R}: dequeue;
    \textsc{W}: enqueue \\
Hilo & \texttt{Thread\{id\}} &
    \textsc{R}: inspeccionar estado;
    \textsc{W}: modificar prioridad/afinidad;
    \textsc{Rv}: destruir \\
IRQ & \texttt{Irq\{line\}} &
    \textsc{R}: esperar interrupción;
    \textsc{W}: acknowledge \\
IoPort & \texttt{IoPort\{base, count\}} &
    \textsc{R}: \texttt{in};
    \textsc{W}: \texttt{out} \\
Notificación & \texttt{Notification\{id\}} &
    \textsc{R}: wait;
    \textsc{W}: signal \\
Dominio & \texttt{Domain\{id\}} &
    \textsc{R}: consultar presupuesto;
    \textsc{W}: añadir/quitar miembros \\
Espacio Addr. & \texttt{AddrSpace\{cr3\}} &
    \textsc{R}: inspeccionar;
    \textsc{W}: map/unmap páginas \\
\bottomrule
\end{tabular}
\end{table}

\section{Invariantes del Sistema de Capacidades}

Enunciamos los invariantes que el sistema de capacidades mantiene en todo momento.

\begin{invariant}[Principio de Menor Privilegio]
\label{inv:least-privilege}
Todo acceso a un objeto del kernel requiere una capacidad válida con derechos suficientes.
No existe un camino de acceso alternativo.
\end{invariant}

\begin{invariant}[No-Amplificación]
\label{inv:no-amplification}
Para toda cadena de delegaciones $c_0 \to c_1 \to \cdots \to c_k$:
$r_k \subseteq r_0$. Los derechos nunca aumentan por delegación.
\end{invariant}

\begin{invariant}[Revocación Total]
\label{inv:total-revocation}
Al revocar $c$, todo descendiente de $c$ en el CDT queda inaccesible. No se produce
una ``fuga'' de capacidades.
\end{invariant}

\begin{invariant}[Consistencia ABA]
\label{inv:aba-consistency}
Si una capacidad $c$ es liberada y su slot es reasignado a $c'$, cualquier handle
obsoleto que apuntaba a $c$ será rechazado (con probabilidad $1 - 2^{-12}$) al
acceder a $c'$.
\end{invariant}


% ====================================================================
\chapter{Comunicación Entre Procesos (IPC)}
\label{chap:ipc}
% ====================================================================

La comunicación entre procesos es la operación más crítica de rendimiento en un
micronúcleo. En un kernel monolítico, los subsistemas se comunican mediante llamadas
a funciones internas; en un micronúcleo, cada interacción entre servidores requiere
un cruce de dominios de protección. La eficiencia del IPC determina, por tanto, la
viabilidad del diseño microkernel.

\sotos{} provee cuatro mecanismos de IPC, cada uno optimizado para un patrón de uso:

\begin{enumerate}
    \item \textbf{Endpoints síncronos}: Mensajes pequeños (72 bytes) transferidos en
          registros durante el cambio de contexto. Cero copias de memoria.
    \item \textbf{Canales asíncronos}: Ring buffers acotados de 16 mensajes para
          transferencia de datos de alto throughput.
    \item \textbf{Notificaciones}: Semáforos binarios para señalización ligera.
    \item \textbf{Mailbox por CPU}: Colas de mensajes entre cores para entrega de IPC
          cross-core con IPI (Inter-Processor Interrupt).
\end{enumerate}

\section{Endpoints Síncronos}

\subsection{Modelo Formal}

\begin{definition}[Endpoint]
\label{def:endpoint}
Un \emph{endpoint} $\epsilon$ es una estructura de sincronización que se encuentra en uno
de tres estados:
\[
    \epsilon.\text{state} \in \{\textsc{Idle}, \textsc{RecvWait}, \textsc{SendWait}\}
\]
con los siguientes campos:
\begin{itemize}
    \item $\epsilon.\text{receiver} \in \{\bot\} \cup \mathbb{T}$: hilo esperando recibir.
    \item $\epsilon.\text{send\_queue} \in \mathbb{T}^*$: cola FIFO de hilos esperando enviar,
          con $|\epsilon.\text{send\_queue}| \leq Q_{\max} = 16$.
    \item $\epsilon.\text{caller} \in \{\bot\} \cup \mathbb{T}$: hilo que ha enviado y
          espera respuesta (protocolo \textsc{Call}).
\end{itemize}
donde $\mathbb{T}$ es el conjunto de identificadores de hilo.
\end{definition}

\begin{definition}[Mensaje IPC]
\label{def:ipc-message}
Un \emph{mensaje} $m = (\text{tag}, \text{regs}, \text{cap})$ consiste en:
\begin{itemize}
    \item $\text{tag} \in \mathbb{N}_{64}$: etiqueta que identifica la operación.
    \item $\text{regs} \in \mathbb{N}_{64}^{8}$: payload de 8 registros (512 bits = 64 bytes).
    \item $\text{cap} \in \{\bot\} \cup \mathbb{N}_{32}$: capacidad opcional a transferir.
\end{itemize}
El tamaño total es $8 + 64 + 4 = 76$ bytes (redondeado a 72 bytes por alineación).
\end{definition}

\subsection{Pools por CPU}

A diferencia de un pool global (que sería un cuello de botella de contención en SMP),
\sotos{} mantiene un pool de endpoints independiente por cada core de CPU.

\begin{definition}[Codificación de Handles de Endpoint]
\label{def:ep-handle}
Un handle de endpoint es un entero de 32 bits:
\[
    h_{\text{ep}} = (\text{core} \ll 28) \bitor (g \ll 20) \bitor s
\]
donde:
\begin{itemize}
    \item $\text{core} = h_{\text{ep}}[31{:}28]$: identificador de CPU (4 bits, máx.\ 16 cores).
    \item $g = h_{\text{ep}}[27{:}20]$: generación (8 bits, wrap-around a 256).
    \item $s = h_{\text{ep}}[19{:}0]$: índice de slot (20 bits, máx.\ 64 por core).
\end{itemize}
\end{definition}

\begin{figure}[ht]
\centering
\begin{tikzpicture}[
    bit/.style={draw, minimum width=0.38cm, minimum height=0.6cm, font=\tiny},
]
% Core bits (31:28)
\foreach \i in {0,...,3} {
    \node[bit, fill=green!20] (core\i) at (\i*0.38, 0) {};
}
% Generation bits (27:20)
\foreach \i in {0,...,7} {
    \node[bit, fill=orange!25] (g\i) at ({(4+\i)*0.38}, 0) {};
}
% Index bits (19:0)
\foreach \i in {0,...,19} {
    \node[bit, fill=blue!15] (s\i) at ({(12+\i)*0.38}, 0) {};
}

\node[font=\small\bfseries, above=3pt] at (core0.north west) {31};
\node[font=\small\bfseries, above=3pt] at (core3.north east) {28};
\node[font=\small\bfseries, above=3pt] at (g0.north west) {27};
\node[font=\small\bfseries, above=3pt] at (g7.north east) {20};
\node[font=\small\bfseries, above=3pt] at (s0.north west) {19};
\node[font=\small\bfseries, above=3pt] at (s19.north east) {0};

\draw[decorate, decoration={brace, amplitude=4pt, mirror}]
    ([yshift=-2pt]core0.south west) -- ([yshift=-2pt]core3.south east)
    node[midway, below=5pt, font=\footnotesize] {CPU};
\draw[decorate, decoration={brace, amplitude=4pt, mirror}]
    ([yshift=-2pt]g0.south west) -- ([yshift=-2pt]g7.south east)
    node[midway, below=5pt, font=\footnotesize] {gen.};
\draw[decorate, decoration={brace, amplitude=4pt, mirror}]
    ([yshift=-2pt]s0.south west) -- ([yshift=-2pt]s19.south east)
    node[midway, below=5pt, font=\footnotesize] {slot};
\end{tikzpicture}
\caption{Codificación del handle de endpoint: 4 bits de core, 8 bits de generación,
20 bits de índice de slot.}
\label{fig:ep-handle}
\end{figure}

\subsection{Semántica de Rendezvous}

El protocolo de endpoints síncronos implementa \emph{rendezvous}: el emisor se bloquea
hasta que un receptor esté listo (y viceversa). La transferencia del mensaje ocurre
directamente entre los registros guardados de ambos hilos --- sin buffer intermedio.

\begin{figure}[ht]
\centering
\begin{tikzpicture}[
    >=Stealth,
    state/.style={draw, rounded corners=5pt, minimum width=2.5cm,
                  minimum height=1cm, font=\small, align=center},
    idle/.style={state, fill=green!15, draw=green!50!black},
    recv/.style={state, fill=blue!15, draw=blue!50!black},
    send/.style={state, fill=orange!15, draw=orange!50!black},
    trans/.style={->, thick, font=\footnotesize},
]

\node[idle] (idle) at (0, 0) {\textsc{Idle}\\{\tiny sin esperas}};
\node[recv] (rw) at (-4, -3) {\textsc{RecvWait}\\{\tiny receptor espera}};
\node[send] (sw) at (4, -3) {\textsc{SendWait}\\{\tiny emisor(es) espera(n)}};

% Idle -> RecvWait
\draw[trans] (idle.south west) -- (rw.north east)
    node[midway, above, sloped] {recv(), sin emisor};

% Idle -> SendWait
\draw[trans] (idle.south east) -- (sw.north west)
    node[midway, above, sloped] {send(), sin receptor};

% RecvWait -> Idle (sender arrives)
\draw[trans, blue!60!black] (rw.north) to[bend left=20]
    node[left, align=right, font=\footnotesize] {send(): rendezvous\\$\to$ wake receptor}
    (idle.west);

% SendWait -> Idle (receiver arrives)
\draw[trans, orange!60!black] (sw.north) to[bend right=20]
    node[right, align=left, font=\footnotesize] {recv(): dequeue emisor\\$\to$ wake emisor}
    (idle.east);

% SendWait -> SendWait (additional senders)
\draw[trans, orange!40!black] (sw.south east) to[bend right=40]
    node[right, font=\footnotesize] {send(): enqueue}
    (sw.south);

% SendWait -> Idle (last sender dequeued)
\path (sw) -- (idle) node[midway, font=\tiny, text=gray] {};

\end{tikzpicture}
\caption{Máquina de estados de un endpoint síncrono. Las transiciones etiquetadas
indican la operación que causa el cambio de estado y la acción resultante.}
\label{fig:ep-state-machine}
\end{figure}

\subsubsection{Operación \textsc{Send}}

\begin{algorithm}[H]
\caption{\textsc{Endpoint-Send}$(\epsilon, m, t_s)$: Enviar mensaje $m$ desde hilo $t_s$.}
\label{alg:ep-send}
\KwIn{Endpoint $\epsilon$, mensaje $m$, hilo emisor $t_s$}
\KwOut{$\textsc{Ok}$ si exitoso, error en caso contrario}
\tcp{Verificar si hay un caller esperando respuesta}
\If{$\epsilon.\text{caller} \neq \bot$}{
    $t_c \gets \epsilon.\text{caller}$; $\epsilon.\text{caller} \gets \bot$\;
    Escribir $m$ en registros de $t_c$\;
    \textsc{Wake}$(t_c)$\;
    \Return $\textsc{Ok}$\tcp*{Respuesta entregada}
}
\Switch{$\epsilon.\text{state}$}{
    \Case{\textsc{RecvWait}}{
        $t_r \gets \epsilon.\text{receiver}$; $\epsilon.\text{receiver} \gets \bot$\;
        $\epsilon.\text{state} \gets \textsc{Idle}$\;
        Escribir $m$ en registros de $t_r$\;
        \textsc{Wake}$(t_r)$\tcp*{Rendezvous: emisor no se bloquea}
    }
    \Case{\textsc{Idle} $|$ \textsc{SendWait}}{
        \lIf{$|\epsilon.\text{send\_queue}| = Q_{\max}$}{\Return $\textsc{Err}(\text{Full})$}
        $\epsilon.\text{send\_queue}.\textsc{Enqueue}(t_s)$\;
        $\epsilon.\text{state} \gets \textsc{SendWait}$\;
        Guardar $m$ en TCB de $t_s$\;
        \textsc{Block}$(t_s)$\tcp*{Emisor se bloquea}
    }
}
\Return $\textsc{Ok}$\;
\end{algorithm}

\subsubsection{Operación \textsc{Recv}}

\begin{algorithm}[H]
\caption{\textsc{Endpoint-Recv}$(\epsilon, t_r)$: Recibir mensaje para hilo $t_r$.}
\label{alg:ep-recv}
\KwIn{Endpoint $\epsilon$, hilo receptor $t_r$}
\KwOut{Mensaje $m$}
\Switch{$\epsilon.\text{state}$}{
    \Case{\textsc{SendWait}}{
        $t_s \gets \epsilon.\text{send\_queue}.\textsc{Dequeue}()$\;
        $\text{is\_caller} \gets (t_s \bitand \text{CALLER\_BIT}) \neq 0$\;
        $t_s \gets t_s \bitand \lnot\text{CALLER\_BIT}$\;
        $m \gets$ leer mensaje guardado en TCB de $t_s$\;
        \eIf{$\text{is\_caller}$}{
            $\epsilon.\text{caller} \gets t_s$\tcp*{Caller espera respuesta}
        }{
            \textsc{Wake}$(t_s)$\tcp*{Send puro: emisor libre}
        }
        \lIf{$|\epsilon.\text{send\_queue}| = 0$}{$\epsilon.\text{state} \gets \textsc{Idle}$}
        \Return $m$\;
    }
    \Case{\textsc{Idle}}{
        $\epsilon.\text{receiver} \gets t_r$\;
        $\epsilon.\text{state} \gets \textsc{RecvWait}$\;
        \textsc{Block}$(t_r)$\tcp*{Receptor se bloquea}
        $m \gets$ leer mensaje de registros de $t_r$\;
        \Return $m$\;
    }
    \Case{\textsc{RecvWait}}{
        \Return $\textsc{Err}(\text{AlreadyWaiting})$\tcp*{Solo un receptor}
    }
}
\end{algorithm}

\subsubsection{Operación \textsc{Call} (RPC)}

\begin{algorithm}[H]
\caption{\textsc{Endpoint-Call}$(\epsilon, m, t_c)$: Enviar $m$ y esperar respuesta.}
\label{alg:ep-call}
\KwIn{Endpoint $\epsilon$, mensaje $m$, hilo caller $t_c$}
\KwOut{Mensaje de respuesta $m_r$}
\Switch{$\epsilon.\text{state}$}{
    \Case{\textsc{RecvWait}}{
        $t_r \gets \epsilon.\text{receiver}$; $\epsilon.\text{receiver} \gets \bot$\;
        $\epsilon.\text{caller} \gets t_c$\;
        $\epsilon.\text{state} \gets \textsc{Idle}$\;
        Escribir $m$ en registros de $t_r$\;
        \textsc{Wake}$(t_r)$\;
        \textsc{Block}$(t_c)$\tcp*{Caller espera respuesta}
    }
    \Case{\textsc{Idle} $|$ \textsc{SendWait}}{
        $\epsilon.\text{send\_queue}.\textsc{Enqueue}(t_c \bitor \text{CALLER\_BIT})$\;
        $\epsilon.\text{state} \gets \textsc{SendWait}$\;
        Guardar $m$ en TCB de $t_c$\;
        \textsc{Block}$(t_c)$\tcp*{Caller espera tanto envío como respuesta}
    }
}
$m_r \gets$ leer respuesta de registros de $t_c$\;
\Return $m_r$\;
\end{algorithm}

\subsection{Propiedades del Protocolo}

\begin{theorem}[Transferencia Sin Copia]
\label{thm:zero-copy}
En el camino rápido (rendezvous inmediato), la operación \textsc{Send} transfiere el
mensaje directamente escribiendo en los registros guardados del receptor. No se realiza
ninguna copia a memoria intermedia.
\end{theorem}

\begin{proof}
Cuando \textsc{Send} encuentra el endpoint en estado \textsc{RecvWait}, la función
$\textsc{write\_ipc\_msg}(t_r, m)$ escribe los campos \texttt{tag} y \texttt{regs} del
mensaje directamente en el campo \texttt{ipc\_msg} del TCB (Thread Control Block) de $t_r$.
Al despertar $t_r$ y ejecutar su retorno de \textsc{Recv}, lee los datos de su propio TCB.

El TCB reside en memoria del kernel (heap del pool de hilos). La escritura es una
asignación de struct --- no una copia de páginas ni una llamada a \texttt{memcpy} de
usuario. La ruta completa es: emisor$\to$TCB-receptor, sin buffer auxiliar.
\end{proof}

\begin{theorem}[Acotación de la Cola de Envío]
\label{thm:bounded-queue}
La cola de envío de cada endpoint tiene cardinalidad máxima $Q_{\max} = 16$. Intentos
adicionales de envío retornan $\textsc{Err}(\text{OutOfResources})$.
\end{theorem}

\begin{proof}
Por inspección del algoritmo \textsc{Endpoint-Send}: antes de encolar, se verifica
$|\epsilon.\text{send\_queue}| < Q_{\max}$. El array subyacente es de tamaño fijo
$\texttt{[Option<ThreadId>; 16]}$.
\end{proof}

\begin{theorem}[Wait-Free en el Camino Rápido]
\label{thm:fast-path}
Si el endpoint está en estado \textsc{RecvWait} al momento de \textsc{Send} (o en
\textsc{SendWait} al momento de \textsc{Recv}), la operación completa en $\bigO(1)$
sin bloquear al invocador.
\end{theorem}

\begin{proof}
En ambos casos, la operación realiza:
\begin{enumerate}
    \item Adquirir el lock del pool por-core: $\bigO(1)$ amortizado (ticket mutex FIFO,
          contención limitada a operaciones en el mismo core).
    \item Leer/escribir campos del endpoint: $\bigO(1)$.
    \item Soltar el lock.
    \item Escribir mensaje en TCB del peer: $\bigO(1)$.
    \item Despertar al peer: $\bigO(1)$ (enqueue en cola de ejecución).
\end{enumerate}
El invocador no se bloquea en ningún punto. La operación es wait-free (no solo lock-free)
para el hilo invocador, ya que completa en un número acotado de pasos sin depender del
progreso de otros hilos.
\end{proof}

\subsection{Mecanismo de Timeout}

\begin{definition}[IPC con Timeout]
\label{def:ipc-timeout}
Una operación \textsc{Call-Timeout}$(\epsilon, m, N)$ establece un plazo absoluto
$\delta = \text{global\_ticks} + N$. Si el hilo sigue bloqueado cuando
$\text{global\_ticks} \geq \delta$, el kernel:
\begin{enumerate}
    \item Marca el flag \texttt{ipc\_timed\_out} en el TCB.
    \item Despierta al hilo.
    \item El hilo detecta el timeout y ejecuta \textsc{Cancel-Caller} para limpiar
          su estado en el endpoint.
\end{enumerate}
\end{definition}

\begin{theorem}[Garantía de Espera Acotada]
\label{thm:bounded-wait}
Para toda invocación de \textsc{Call-Timeout}$(\epsilon, m, N)$ con $N > 0$, el hilo
invocador se desbloquea en a lo sumo $N$ ticks del temporizador global, independientemente
del comportamiento del servidor.
\end{theorem}

\begin{proof}
Sea $t_0$ el tick global al momento de la invocación. Se establece $\delta = t_0 + N$.
El temporizador LAPIC genera interrupciones periódicas que incrementan
\texttt{global\_ticks} atómicamente. En la rutina del tick del planificador:

\begin{enumerate}
    \item Para cada hilo bloqueado con $\texttt{ipc\_timeout} > 0$:
    \item Si $\texttt{global\_ticks} \geq \texttt{ipc\_timeout}$:
    \item Marcar $\texttt{ipc\_timed\_out} \gets \texttt{true}$.
    \item Llamar $\textsc{Wake}(t)$.
\end{enumerate}

Como \texttt{global\_ticks} es monótonamente creciente (AtomicU64 con fetch\_add) y el
timer tick ocurre con frecuencia $f$ (100~Hz por defecto), el tiempo de pared hasta el
timeout es $\leq N/f + 1/f$ segundos (un tick de latencia adicional por granularidad).
El hilo se desbloquea garantizadamente.
\end{proof}

\section{Canales Asíncronos}

\begin{definition}[Canal]
\label{def:channel}
Un \emph{canal} es un ring buffer acotado de capacidad $K = 16$ mensajes, con semántica
producteur-consommateur:
\begin{itemize}
    \item $\textsc{Ch-Send}$: encola un mensaje. Bloquea si está lleno.
    \item $\textsc{Ch-Recv}$: desencola un mensaje. Bloquea si está vacío.
\end{itemize}
\end{definition}

\begin{proposition}[Propiedades del Canal]
\label{prop:channel-props}
Sea $C$ un canal con capacidad $K$. Entonces:
\begin{enumerate}
    \item $0 \leq |C| \leq K$ en todo momento (acotación).
    \item Los mensajes se entregan en orden FIFO (ring buffer circular).
    \item En ausencia de bloqueo, \textsc{Ch-Send} y \textsc{Ch-Recv} son $\bigO(1)$.
\end{enumerate}
\end{proposition}

\begin{proof}
\leavevmode
\begin{enumerate}
    \item La cola empieza con $|C| = 0$. \textsc{Ch-Send} incrementa $|C|$ solo si
          $|C| < K$ (bloquea en caso contrario). \textsc{Ch-Recv} decrementa $|C|$ solo
          si $|C| > 0$ (bloquea en caso contrario). Ambas operaciones están protegidas
          por un mutex, por lo que no hay carreras.
    \item El ring buffer usa índices \texttt{head} y \texttt{tail} con aritmética modular:
          $\texttt{tail} = (\texttt{tail} + 1) \bmod K$ y
          $\texttt{head} = (\texttt{head} + 1) \bmod K$. El consumidor siempre lee
          desde \texttt{head} (el mensaje más antiguo).
    \item Acceso indexado al array $\texttt{buffer}[i]$ es $\bigO(1)$. La aritmética
          modular es $\bigO(1)$.
\end{enumerate}
\end{proof}

\section{Transferencia de Capacidades por IPC}

Una característica fundamental del sistema es la capacidad de transferir capacidades
entre procesos a través de mensajes IPC.

\begin{definition}[Transferencia de Capacidad]
\label{def:cap-transfer}
Cuando un mensaje $m$ incluye $m.\text{cap} = c_s$ (handle de una capacidad fuente), el
kernel automáticamente:
\begin{enumerate}
    \item Verifica que $c_s$ tenga el derecho \textsc{Grant}.
    \item Crea una nueva capacidad $c_d = \textsc{Grant}(c_s, \rightsall)$ como hija de $c_s$.
    \item Reemplaza $m.\text{cap}$ con el handle de $c_d$.
\end{enumerate}
El receptor obtiene una capacidad derivada completa (todos los derechos del padre).
\end{definition}

\begin{proposition}[Seguridad de la Transferencia]
\label{prop:transfer-security}
La transferencia de capacidades por IPC preserva todos los invariantes del sistema:
\begin{enumerate}
    \item \emph{No-amplificación}: $c_d$ tiene derechos $\leq c_s$ (por \cref{lem:rights-monotone}).
    \item \emph{Revocabilidad}: $c_d$ es hija de $c_s$ en el CDT, por lo que revocar $c_s$
          revocará $c_d$ (por \cref{thm:revoke-correct}).
    \item \emph{Autenticidad}: Solo el poseedor de $c_s$ (con derecho \textsc{Grant})
          puede iniciar la transferencia.
\end{enumerate}
\end{proposition}

\section{Resumen de Complejidades}

\begin{table}[ht]
\centering
\caption{Complejidad temporal de las operaciones de IPC.}
\label{tab:ipc-complexity}
\begin{tabular}{@{}llll@{}}
\toprule
\textbf{Operación} & \textbf{Camino rápido} & \textbf{Camino lento} & \textbf{Notas} \\
\midrule
\textsc{Send} (rendezvous) & $\bigO(1)$ & $\bigO(1)$ + bloqueo & Wait-free si receptor listo \\
\textsc{Recv} (rendezvous) & $\bigO(1)$ & $\bigO(1)$ + bloqueo & Wait-free si emisor listo \\
\textsc{Call} (RPC) & $\bigO(1)$ & $\bigO(1)$ + 2 bloqueos & Send + esperar respuesta \\
\textsc{Call-Timeout} & $\bigO(1)$ & $\leq N$ ticks & Plazo garantizado \\
\textsc{Ch-Send} & $\bigO(1)$ & $\bigO(1)$ + bloqueo & Buffer no lleno \\
\textsc{Ch-Recv} & $\bigO(1)$ & $\bigO(1)$ + bloqueo & Buffer no vacío \\
Cap transfer & $\bigO(1)$ & --- & Verificación + inserción \\
\bottomrule
\end{tabular}
\end{table}


% ====================================================================
\chapter{Planificación}
\label{chap:scheduling}
% ====================================================================

El planificador de \sotos{} gestiona la multiplexación temporal del (o los) procesador(es)
entre hilos concurrentes. Dado que todos los servicios del sistema operativo se ejecutan
como procesos de usuario, el planificador es el árbitro final de los recursos de cómputo.

\section{Multi-Level Queue (MLQ)}

\subsection{Estructura de Colas}

\begin{definition}[Cola Multinivel]
\label{def:mlq}
El planificador mantiene, por cada CPU, un arreglo de $P = 4$ colas FIFO independientes
indexadas por \emph{clase de prioridad}:
\[
    Q_{\text{cpu}}[k], \quad k \in \{0, 1, 2, 3\}
\]
donde las clases son:
\begin{center}
\begin{tabular}{cll}
\toprule
$k$ & \textbf{Clase} & \textbf{Rango de prioridad} \\
\midrule
0 & Tiempo real  & $[0, 63]$ \\
1 & Alta         & $[64, 127]$ \\
2 & Normal       & $[128, 191]$ \\
3 & Baja         & $[192, 255]$ \\
\bottomrule
\end{tabular}
\end{center}
\end{definition}

La función de mapeo es:
\[
    \text{class}(p) = \min\left(\left\lfloor \frac{p}{64} \right\rfloor, 3\right)
\]

\begin{figure}[ht]
\centering
\begin{tikzpicture}[
    >=Stealth,
    queue/.style={draw, minimum width=5cm, minimum height=0.7cm, font=\small},
    thread/.style={draw, fill=blue!15, minimum width=0.8cm, minimum height=0.5cm,
                   font=\tiny},
    label/.style={font=\small\bfseries, anchor=east},
    steal/.style={->, thick, red!60!black, dashed, bend left=20},
]

% CPU 0
\node[font=\bfseries] at (3, 4) {CPU 0};
\node[label] at (0, 3) {RT};
\node[queue, fill=red!8] (q00) at (3, 3) {};
\node[label] at (0, 2.1) {Alta};
\node[queue, fill=orange!8] (q01) at (3, 2.1) {};
\node[label] at (0, 1.2) {Normal};
\node[queue, fill=yellow!12] (q02) at (3, 1.2) {};
\node[label] at (0, 0.3) {Baja};
\node[queue, fill=green!8] (q03) at (3, 0.3) {};

% Threads in CPU 0
\node[thread, fill=red!25] at (1.2, 3) {$t_1$};
\node[thread, fill=red!25] at (2.0, 3) {$t_2$};
\node[thread, fill=orange!25] at (1.2, 2.1) {$t_3$};
\node[thread, fill=yellow!30] at (1.2, 1.2) {$t_4$};
\node[thread, fill=yellow!30] at (2.0, 1.2) {$t_5$};
\node[thread, fill=yellow!30] at (2.8, 1.2) {$t_6$};
\node[thread, fill=green!25] at (1.2, 0.3) {$t_7$};

% CPU 1
\node[font=\bfseries] at (10, 4) {CPU 1};
\node[queue, fill=red!8] (q10) at (10, 3) {};
\node[queue, fill=orange!8] (q11) at (10, 2.1) {};
\node[queue, fill=yellow!12] (q12) at (10, 1.2) {};
\node[queue, fill=green!8] (q13) at (10, 0.3) {};

% Threads in CPU 1 (fewer)
\node[thread, fill=yellow!30] at (8.2, 1.2) {$t_8$};

% Labels for CPU 1
\node[label] at (7, 3) {RT};
\node[label] at (7, 2.1) {Alta};
\node[label] at (7, 1.2) {Normal};
\node[label] at (7, 0.3) {Baja};

% Work stealing arrow
\draw[steal] (q03.east) to
    node[above, font=\footnotesize, text=red!60!black] {robo de trabajo}
    (q13.west);

% Dequeue arrow
\draw[->, thick, blue!60!black] (q00.west) -- ++(-0.8, 0)
    node[left, font=\footnotesize] {dequeue};

% Annotation
\node[font=\footnotesize\itshape, text=gray, anchor=north] at (6.5, -0.5) {
    CPU 1 roba de la cola de menor prioridad de CPU 0
};

\end{tikzpicture}
\caption{Colas multinivel por CPU con robo de trabajo. Cada CPU despacha primero de su
cola de mayor prioridad (RT). Cuando la CPU queda ociosa, roba del extremo de menor
prioridad de otra CPU.}
\label{fig:mlq}
\end{figure}

\subsection{Algoritmo de Despacho}

\begin{algorithm}[H]
\caption{\textsc{Dequeue-Local}$(Q_{\text{cpu}})$: Despachar el siguiente hilo de la CPU local.}
\label{alg:dequeue-local}
\KwIn{Colas $Q_{\text{cpu}}[0..3]$ de la CPU actual}
\KwOut{Índice de hilo o $\bot$}
\For{$k \gets 0$ \KwTo $3$}{
    \If{$Q_{\text{cpu}}[k] \neq \emptyset$}{
        \Return $Q_{\text{cpu}}[k].\textsc{PopFront}()$\;
    }
}
\Return $\bot$\;
\end{algorithm}

El algoritmo siempre despacha del nivel de mayor prioridad, lo que proporciona preemption
estricta: un hilo de tiempo real siempre se ejecuta antes que uno normal.

\section{Robo de Trabajo (Work Stealing)}

\subsection{Algoritmo}

Cuando una CPU agota su cola local, intenta ``robar'' trabajo de otras CPUs.

\begin{algorithm}[H]
\caption{\textsc{Dequeue-Any}$(\text{my\_cpu})$: Despachar con robo de trabajo.}
\label{alg:dequeue-any}
\KwIn{Índice de la CPU actual $\text{my\_cpu}$}
\KwOut{Índice de hilo o $\bot$}
\tcp{Intentar cola local primero}
$t \gets \textsc{Dequeue-Local}(Q_{\text{my\_cpu}})$\;
\lIf{$t \neq \bot$}{\Return $t$}
\tcp{Robar de otras CPUs (round-robin)}
\For{$\text{offset} \gets 1$ \KwTo $C - 1$}{
    $\text{victim} \gets (\text{my\_cpu} + \text{offset}) \bmod C$\;
    \If{$\textsc{TryLock}(Q_{\text{victim}})$ tiene éxito}{
        $t \gets Q_{\text{victim}}.\textsc{PopBackLowest}()$\;
        \textsc{Unlock}$(Q_{\text{victim}})$\;
        \lIf{$t \neq \bot$}{\Return $t$}
    }
}
\Return $\bot$\tcp*{Todas las CPUs ociosas}
\end{algorithm}

Aspectos clave del diseño:

\begin{enumerate}
    \item \textbf{Selección de víctima}: Round-robin simple ($\text{my\_cpu} + 1, +2, \ldots$),
          lo que distribuye la carga de robo uniformemente.
    \item \textbf{Selección de tarea}: Se roba del \emph{extremo trasero} de la cola de
          \emph{menor prioridad} (\textsc{PopBackLowest}). Esto tiene dos beneficios:
          (a) las tareas de alta prioridad de la víctima no se perturban; y
          (b) las tareas antiguas (``frías'' en cache de la víctima) tienen menor costo
          de migración.
    \item \textbf{Try-lock}: Se usa \textsc{TryLock} (no bloqueo) para evitar inversión
          de prioridad y deadlock si la víctima también está operando sobre su cola.
\end{enumerate}

\subsection{Análisis de Rendimiento}

\begin{theorem}[Cota de Makespan con Robo de Trabajo]
\label{thm:work-stealing}
Sea un programa con trabajo total $T_1$ (suma de tiempos de todas las tareas) y
profundidad $T_\infty$ (longitud del camino crítico). Con $P$ procesadores y robo de
trabajo aleatorizado, el makespan esperado satisface:
\[
    \mathbb{E}[T_P] = \bigO\!\left(\frac{T_1}{P} + T_\infty\right).
\]
\end{theorem}

\begin{proof}[Referencia de la demostración]
Este resultado fue demostrado por Blumofe y Leiserson~\cite{blumofe1999} para el
planificador de Cilk-5. La prueba se basa en una \emph{analogía de bola-en-urna}:

\begin{enumerate}
    \item En cada paso temporal, si una CPU tiene trabajo local, lo ejecuta (paso productivo).
    \item Si está ociosa, roba de una CPU elegida uniformemente al azar.
    \item Un argumento de potencial muestra que el número esperado de robos es
          $\bigO(P \cdot T_\infty)$.
    \item Cada paso es o productivo o un intento de robo. Los pasos productivos suman
          $T_1$ (el trabajo total). Los robos suman $\bigO(P \cdot T_\infty)$ en expectativa.
    \item El makespan total es $(T_1 + \bigO(P \cdot T_\infty)) / P =
          \bigO(T_1/P + T_\infty)$.
\end{enumerate}

\sotos{} usa selección round-robin en lugar de aleatoria, lo que provee propiedades
similares en la práctica para $P$ pequeño (típicamente $\leq 16$ cores).
\end{proof}

\begin{remark}
La cota $\bigO(T_1/P + T_\infty)$ es óptima en el sentido de que:
\begin{itemize}
    \item $T_1/P$ es el tiempo mínimo con paralelismo perfecto (trabajo dividido equitativamente).
    \item $T_\infty$ es el tiempo mínimo con infinitos procesadores (solo el camino crítico).
\end{itemize}
El speedup lineal se alcanza cuando $T_1/T_\infty \gg P$ (el \emph{paralelismo}
$T_1/T_\infty$ excede el número de procesadores).
\end{remark}

\section{Cambio de Contexto}

\subsection{Salvar y Restaurar Registros}

El cambio de contexto es una operación a nivel de ensamblador que salva los registros
callee-saved del hilo actual y restaura los del hilo destino.

\begin{definition}[TrapFrame]
\label{def:trapframe}
El \emph{TrapFrame} almacenado en la pila del kernel durante un cambio de contexto contiene
los 6 registros callee-saved del ABI System V x86-64:
\[
    \text{TrapFrame} = (\texttt{r15}, \texttt{r14}, \texttt{r13}, \texttt{r12},
                        \texttt{rbx}, \texttt{rbp})
\]
más la dirección de retorno implícita en el tope de la pila.
\end{definition}

\begin{lstlisting}[language={[x86masm]Assembler}, caption={Rutina de cambio de contexto
(\texttt{context\_switch}). RDI = puntero a RSP antiguo; RSI = RSP nuevo.},
label={lst:context-switch}]
context_switch:
    push rbp
    push rbx
    push r12
    push r13
    push r14
    push r15
    mov [rdi], rsp      ; salvar RSP del hilo actual
    mov rsp, rsi        ; cargar RSP del hilo destino
    pop r15
    pop r14
    pop r13
    pop r12
    pop rbx
    pop rbp
    ret                  ; saltar a la continuación del nuevo hilo
\end{lstlisting}

\begin{proposition}[Costo del Cambio de Contexto]
\label{prop:ctx-switch-cost}
El cambio de contexto de \sotos{} requiere exactamente:
\begin{itemize}
    \item 6 instrucciones \texttt{push} + 1 \texttt{mov [mem], reg}: salvar contexto (7 uops).
    \item 1 \texttt{mov reg, mem} + 6 instrucciones \texttt{pop} + 1 \texttt{ret}:
          restaurar contexto (8 uops).
    \item Total: 15 uops $\approx$ 5--10 nanosegundos en hardware moderno.
\end{itemize}
\end{proposition}

\subsection{Transición a Modo Usuario: SYSCALL/SYSRET}

Para hilos de usuario, el cambio de contexto termina con una transición de Ring~0 a Ring~3
mediante la instrucción \texttt{sysretq}.

\begin{definition}[Configuración de MSRs para SYSCALL/SYSRET]
\label{def:syscall-msr}
Las siguientes Model-Specific Registers (MSRs) se configuran durante la inicialización:
\begin{itemize}
    \item \texttt{IA32\_STAR} (0xC0000081): Segmentos para SYSCALL (kernel CS/SS) y
          SYSRET (user CS/SS).
    \item \texttt{IA32\_LSTAR} (0xC0000082): Dirección del punto de entrada del kernel
          para SYSCALL.
    \item \texttt{IA32\_FMASK} (0xC0000084): Máscara de RFLAGS aplicada en SYSCALL
          (deshabilita interrupciones durante la transición).
\end{itemize}
\end{definition}

El flujo completo de un \texttt{sysretq} es:

\begin{enumerate}
    \item Cargar CR3 del hilo destino (cambio de espacio de direcciones).
    \item Cargar \texttt{rcx} $\gets$ RIP del usuario.
    \item Cargar \texttt{r11} $\gets$ \texttt{0x202} (RFLAGS con IF=1, habilita interrupciones).
    \item Cargar \texttt{rsp} $\gets$ RSP del usuario.
    \item Ejecutar \texttt{swapgs} (intercambiar GS base kernel/usuario).
    \item Ejecutar \texttt{sysretq} (transición a Ring~3).
\end{enumerate}

\subsection{Restauración de TLS y Estado Extendido}

Además de los registros de propósito general, el cambio de contexto restaura:

\begin{itemize}
    \item \texttt{IA32\_FS\_BASE} (MSR 0xC0000100): Base del segmento FS, usado por
          Thread-Local Storage (TLS). Cada hilo tiene su propio valor, restaurado via
          \texttt{wrmsr}.
    \item \texttt{IA32\_KERNEL\_GS\_BASE} (MSR 0xC0000102): Valor de GS del usuario
          (guardado mientras el kernel usa GS para \texttt{PerCpu}).
    \item Estado FPU/SSE: 512 bytes salvados con \texttt{fxsave64} y restaurados con
          \texttt{fxrstor64}, necesarios porque instrucciones SIMD como \texttt{movaps}
          (usadas por musl/Rust para \texttt{memcpy}) corrompen registros XMM entre
          cambios de contexto.
\end{itemize}

\section{Dominios de Planificación}

\begin{definition}[Dominio de Planificación]
\label{def:sched-domain}
Un \emph{dominio de planificación} $D = (Q_D, P_D, M_D)$ donde:
\begin{itemize}
    \item $Q_D \in \mathbb{N}$: quantum de ticks por período.
    \item $P_D \in \mathbb{N}$: longitud del período en ticks.
    \item $M_D \subseteq \mathbb{T}$: conjunto de hilos miembros ($|M_D| \leq 64$).
\end{itemize}
El dominio mantiene un contador $\text{consumed} \leq Q_D$ de ticks consumidos en el
período actual. Cuando $\text{consumed} = Q_D$, el dominio entra en estado
\textsc{Depleted} y todos sus hilos son suspendidos hasta el siguiente refill
(cuando $\text{global\_ticks} \geq \text{next\_refill}$).
\end{definition}

\begin{theorem}[Aislamiento Temporal entre Dominios]
\label{thm:domain-isolation}
Sea $D_1, D_2$ dos dominios con cuotas $Q_1, Q_2$ y período $P$. Entonces en cualquier
ventana de $P$ ticks, los hilos de $D_i$ consumen a lo sumo $Q_i$ ticks de CPU.
\end{theorem}

\begin{proof}
El planificador incrementa $D.\text{consumed}$ en cada tick donde un hilo de $D$ se está
ejecutando. Cuando $D.\text{consumed} = Q_D$, todos los hilos de $D$ son suspendidos
(movidos a $D.\text{suspended}$) y no son elegibles para despacho (la función
\textsc{Dequeue-Non-Depleted} los filtra). Por lo tanto, en cada período de $P$ ticks,
ningún hilo de $D$ puede ejecutar más de $Q_D$ ticks.

La restauración ocurre cuando $\text{global\_ticks} \geq \text{next\_refill}$:
$D.\text{consumed} \gets 0$, $D.\text{state} \gets \textsc{Active}$,
$D.\text{next\_refill} \gets \text{next\_refill} + P$, y todos los hilos suspendidos
son re-encolados.
\end{proof}

\section{Límites de Recursos por Hilo}

Cada TCB mantiene contadores de consumo individual:

\begin{definition}[Límites de Recursos]
\label{def:resource-limits}
Para cada hilo $t$:
\begin{itemize}
    \item $t.\texttt{cpu\_ticks}$: ticks de CPU consumidos (acumulativo).
    \item $t.\texttt{cpu\_tick\_limit}$: límite máximo (0 = ilimitado).
    \item $t.\texttt{mem\_pages}$: páginas de memoria asignadas.
    \item $t.\texttt{mem\_page\_limit}$: límite máximo (0 = ilimitado).
\end{itemize}
\end{definition}

\begin{theorem}[El Agotamiento de Recursos Termina al Hilo Infractor]
\label{thm:resource-exhaustion}
Si un hilo $t$ alcanza $t.\texttt{cpu\_ticks} \geq t.\texttt{cpu\_tick\_limit} > 0$,
el planificador lo marca como \textsc{Dead}. Ningún otro hilo se ve afectado.
\end{theorem}

\begin{proof}
En la función \textsc{Tick} del planificador (invocada en cada interrupción de timer):

\begin{enumerate}
    \item Se incrementa $t.\texttt{cpu\_ticks}$ para el hilo actual.
    \item Se verifica: si $t.\texttt{cpu\_tick\_limit} > 0$ y
          $t.\texttt{cpu\_ticks} \geq t.\texttt{cpu\_tick\_limit}$, entonces
          $t.\texttt{state} \gets \textsc{Dead}$.
    \item El hilo muerto nunca es re-encolado (la función \textsc{Enqueue} no acepta
          hilos en estado \textsc{Dead}).
\end{enumerate}

Los demás hilos no son afectados porque: (a) cada hilo tiene contadores independientes; y
(b) el planificador preemptivo cambia al siguiente hilo listo inmediatamente después de
marcar al infractor como muerto.
\end{proof}

\section{Asignación de Marcos Físicos}
\label{sec:frame-alloc}

Aunque la gestión de memoria virtual (paginación) se delega al servidor VMM en espacio
de usuario, la asignación de \emph{marcos físicos} --- las unidades atómicas de memoria
física de 4\,KiB --- es una primitiva del kernel.

\subsection{Asignador de Bitmap}

\begin{definition}[Bitmap de Marcos]
\label{def:frame-bitmap}
El asignador mantiene un arreglo de bits $B[0..N)$ donde $N$ es el número total de marcos
físicos en el sistema:
\[
    B[i] = \begin{cases}
        0 & \text{si el marco } i \text{ está libre} \\
        1 & \text{si el marco } i \text{ está asignado}
    \end{cases}
\]
Un marco tiene dirección física $\text{addr}(i) = i \times 4096$.
\end{definition}

\begin{algorithm}[H]
\caption{\textsc{Frame-Alloc}$(B)$: Asignar un marco físico libre.}
\label{alg:frame-alloc}
\KwIn{Bitmap $B$ con hint $h$}
\KwOut{Marco físico o $\bot$}
\tcp{Dos pasadas: hint..N, luego 0..hint}
\For{$\text{pass} \gets 0$ \KwTo $1$}{
    $(s, e) \gets$ \textbf{if} pass $= 0$ \textbf{then} $(h, N)$ \textbf{else} $(0, h)$\;
    $i \gets s$\;
    \While{$i < e$}{
        $\text{byte} \gets i / 8$\;
        \If{$B[\text{byte}] = \texttt{0xFF}$}{
            $i \gets (\text{byte} + 1) \times 8$\tcp*{Skip: byte completo ocupado}
            \textbf{continue}\;
        }
        \For{$b \gets i \bmod 8$ \KwTo $7$}{
            $j \gets \text{byte} \times 8 + b$\;
            \If{$j \geq e$}{\textbf{break}}
            \If{$B[j] = 0$}{
                $B[j] \gets 1$; $h \gets j + 1$\;
                \Return $\text{PhysFrame}(j \times 4096)$\;
            }
        }
        $i \gets (\text{byte} + 1) \times 8$\;
    }
}
\Return $\bot$\tcp*{Sin marcos libres}
\end{algorithm}

\begin{proposition}[Complejidad del Asignador de Marcos]
\label{prop:frame-alloc-complexity}
\leavevmode
\begin{enumerate}
    \item \textsc{Frame-Alloc}: $\bigO(N/8)$ peor caso (escaneo completo del bitmap a
          nivel de byte), $\bigO(1)$ caso típico (el hint apunta a un marco libre).
    \item \textsc{Frame-Free}: $\bigO(1)$ (acceso directo por índice).
    \item \textsc{Frame-Alloc-Contiguous}$(k)$: $\bigO(N \cdot k)$ peor caso
          (buscar secuencia de $k$ marcos libres consecutivos).
\end{enumerate}
\end{proposition}

\begin{proof}
\leavevmode
\begin{enumerate}
    \item La optimización de \emph{byte skip} ($B[\text{byte}] = \texttt{0xFF}$) reduce
          el escaneo a $N/8$ comparaciones en el peor caso. El hint $h$ memoriza la
          posición del último marco asignado, proporcionando localidad temporal: si
          hay marcos libres cerca del último asignado, la búsqueda es $\bigO(1)$.
    \item \textsc{Frame-Free} calcula $\text{byte} = i/8$, $\text{bit} = i \bmod 8$ y
          limpia $B[\text{byte}][\text{bit}]$: una división, un AND y una escritura.
    \item La búsqueda contigua verifica cada posición candidata $i$ y luego escanea
          $k$ bits consecutivos. En el peor caso (bitmap casi lleno), cada candidato
          falla y se necesitan $N - k + 1$ intentos de $k$ verificaciones cada uno.
\end{enumerate}
\end{proof}

\subsection{Inicialización desde el Mapa de Memoria}

El asignador se inicializa a partir del mapa de memoria proporcionado por el cargador
Limine:

\begin{enumerate}
    \item Se determina la dirección física más alta de RAM utilizable para dimensionar
          el bitmap: $N = \lceil \text{max\_addr} / 4096 \rceil$.
    \item Se asigna el bitmap ($\lceil N/8 \rceil$ bytes) dentro de una región utilizable.
    \item Se marca \emph{todo} como ocupado (\texttt{memset(0xFF)}).
    \item Se recorren las regiones \texttt{USABLE} del mapa de memoria, liberando los
          marcos correspondientes (excepto los que contienen el bitmap mismo).
\end{enumerate}

\begin{invariant}[Consistencia del Bitmap]
\label{inv:bitmap-consistency}
En todo momento, $B[i] = 1$ si y solo si el marco $i$ ha sido asignado y no liberado.
Un marco asignado por \textsc{Frame-Alloc} no puede ser asignado nuevamente hasta que
sea liberado por \textsc{Frame-Free}.
\end{invariant}

\begin{proof}
Por inducción sobre la secuencia de operaciones:
\begin{itemize}
    \item \textbf{Base}: Después de la inicialización, todos los marcos de regiones
          no utilizables tienen $B[i] = 1$ (marcados como ocupados). Los marcos de
          regiones utilizables tienen $B[i] = 0$ (libres). El bitmap del asignador
          mismo está marcado como ocupado.
    \item \textbf{Paso (Alloc)}: \textsc{Frame-Alloc} solo selecciona $i$ con $B[i] = 0$,
          luego establece $B[i] = 1$. No se cambia ningún otro bit.
    \item \textbf{Paso (Free)}: \textsc{Frame-Free} verifica que $B[i] = 1$ (assertion
          de doble liberación) y establece $B[i] = 0$. No se cambia ningún otro bit.
\end{itemize}
Todas las operaciones están protegidas por un spinlock global, por lo que no hay carreras
de datos.
\end{proof}

\section{Resumen del Capítulo}

Este capítulo ha presentado los cuatro mecanismos centrales del planificador de \sotos{}:

\begin{enumerate}
    \item \textbf{MLQ por CPU}: Cuatro niveles de prioridad con despacho estricto por
          prioridad, proporcionando determinismo para hilos de tiempo real.
    \item \textbf{Robo de trabajo}: Balanceo de carga automático con cota teórica
          $\bigO(T_1/P + T_\infty)$.
    \item \textbf{Dominios de planificación}: Aislamiento temporal con cuotas periódicas,
          garantizando que un dominio no puede monopolizar la CPU.
    \item \textbf{Límites de recursos}: Protección individual por hilo contra consumo
          excesivo de CPU y memoria.
\end{enumerate}

Junto con el asignador de marcos físicos, estos mecanismos completan las cinco primitivas
del kernel descritas en el \cref{chap:intro}.

% ####################################################################
% ####################################################################
%
%              PARTE II — ESPACIO DE USUARIO
%
% ####################################################################
% ####################################################################

\part{Espacio de Usuario}

% ====================================================================
\chapter{Gestión de Memoria}
\label{chap:memoria}
% ====================================================================

La gestión de memoria en \sotos{} se divide en dos niveles:
\begin{enumerate}
    \item \textbf{Asignador de marcos físicos} (en el núcleo): administra las páginas
          de 4\,KiB de RAM física mediante un bitmap.
    \item \textbf{Tablas de páginas y memoria virtual} (en espacio de usuario): el
          servidor VMM maneja los fallos de página, la paginación bajo demanda y el
          Copy-on-Write.
\end{enumerate}

Esta separación de competencias es la decisión arquitectónica central del micronúcleo:
el kernel nunca interpreta direcciones virtuales de usuario, delegando toda la política
de memoria virtual al servidor VMM.

% --------------------------------------------------------------------
\section{Asignador de Marcos Físicos}
\label{sec:frame-alloc}

\subsection{Estructura de Datos}

El asignador utiliza un bitmap global protegido por un \texttt{spin::Mutex}.
Cada bit representa un marco de 4\,KiB: bit $0$ indica libre, bit $1$ indica
ocupado.

\begin{definition}[Bitmap de marcos]
Sea $N = \lceil \textit{max\_addr} / 4096 \rceil$ el número total de marcos
rastreables. El bitmap $B[0..N-1]$ es un arreglo de $\lceil N/8 \rceil$ bytes,
donde el bit $j$ del byte $i$ corresponde al marco $8i + j$.
\end{definition}

Además del bitmap, el asignador mantiene:
\begin{itemize}
    \item \texttt{total\_frames}: número total de marcos rastreados.
    \item \texttt{free\_frames}: contador de marcos libres (decrementado en alloc, incrementado en free).
    \item \texttt{search\_hint}: índice del último marco asignado + 1, punto de inicio para la siguiente búsqueda.
\end{itemize}

\subsection{Algoritmo de Asignación}

\begin{algorithm}[H]
\caption{\textsc{Frame-Alloc}: Asignación de un marco físico}
\label{alg:frame-alloc}
\KwIn{Bitmap $B[0..N-1]$, puntero de búsqueda \texttt{hint}}
\KwOut{Dirección física del marco asignado, o error}
\For{$\text{pass} \leftarrow 0$ \KwTo $1$}{
    $(s, e) \leftarrow \begin{cases} (\texttt{hint}, N) & \text{si pass} = 0 \\ (0, \texttt{hint}) & \text{si pass} = 1 \end{cases}$ \\
    $i \leftarrow s$ \\
    \While{$i < e$}{
        $\texttt{byte} \leftarrow \lfloor i / 8 \rfloor$ \\
        \If{$B[\texttt{byte}] = \texttt{0xFF}$}{
            $i \leftarrow (\texttt{byte} + 1) \times 8$ \tcp*{saltar byte completo}
            \textbf{continuar} \\
        }
        \For{$b \leftarrow (i \bmod 8)$ \KwTo $7$}{
            $\texttt{idx} \leftarrow \texttt{byte} \times 8 + b$ \\
            \If{$\texttt{idx} \geq e$}{\textbf{salir del bucle}}
            \If{$\text{bit}(B, \texttt{idx}) = 0$}{
                $\text{set\_bit}(B, \texttt{idx})$ \\
                $\texttt{free\_frames} \leftarrow \texttt{free\_frames} - 1$ \\
                $\texttt{hint} \leftarrow \texttt{idx} + 1$ \\
                \Return $\texttt{idx} \times 4096$
            }
        }
        $i \leftarrow (\texttt{byte} + 1) \times 8$
    }
}
\Return \texttt{None} \tcp*{sin marcos libres}
\end{algorithm}

\begin{theorem}[Complejidad del asignador bitmap]
\label{thm:bitmap-complexity}
Sea $N$ el número total de marcos y $F$ el número de marcos libres.
\begin{enumerate}
    \item En el peor caso, \textsc{Frame-Alloc} es $\bigO(N/8)$ (escaneo completo,
          byte a byte).
    \item En el caso amortizado con puntero de búsqueda (\texttt{hint}),
          si las liberaciones son infrecuentes respecto a las asignaciones,
          \textsc{Frame-Alloc} es $\bigO(1)$ amortizado.
\end{enumerate}
\end{theorem}

\begin{proof}
\begin{enumerate}
    \item \textbf{Peor caso.} El bucle exterior itera sobre como máximo $\lceil N/8 \rceil$ bytes
          (dos pasadas concatenadas cubren todos los bytes una vez).
          Cada byte se examina en tiempo $\bigO(1)$: si vale \texttt{0xFF} se salta;
          si no, se busca el primer bit libre dentro de 8 posiciones.
          Total: $\bigO(N/8)$.

    \item \textbf{Amortizado.} Consideremos una secuencia de $k$ asignaciones consecutivas
          sin liberaciones. Después de cada asignación, \texttt{hint} se adelanta al
          siguiente índice. La siguiente asignación comienza en \texttt{hint} y encuentra
          un marco libre en $\bigO(1)$ si los marcos están contiguos (caso común en un
          bitmap recién inicializado). La suma total de trabajo para $k$ asignaciones
          es a lo sumo $N/8$ (un recorrido completo), dando $\bigO(N/(8k))$ amortizado
          por operación. Cuando $k \geq N/8$, el costo amortizado es $\bigO(1)$. \qedhere
\end{enumerate}
\end{proof}

\subsection{Asignación Contigua para DMA}

Ciertos dispositivos (virtqueues de virtio, PRP lists de NVMe) requieren marcos
físicamente contiguos. La función \texttt{alloc\_contiguous(count)} extiende el
algoritmo anterior:

\begin{algorithm}[H]
\caption{\textsc{Frame-Alloc-Contiguous}: Asignación de $k$ marcos contiguos}
\label{alg:frame-alloc-contig}
\KwIn{Bitmap $B[0..N-1]$, cantidad $k$}
\KwOut{Dirección física base, o error}
$\texttt{run} \leftarrow 0$ \\
$\texttt{run\_start} \leftarrow 0$ \\
\For{$i \leftarrow 0$ \KwTo $N-1$}{
    \eIf{$\text{bit}(B, i) = 0$}{
        \If{$\texttt{run} = 0$}{$\texttt{run\_start} \leftarrow i$}
        $\texttt{run} \leftarrow \texttt{run} + 1$ \\
        \If{$\texttt{run} = k$}{
            \For{$j \leftarrow \texttt{run\_start}$ \KwTo $\texttt{run\_start} + k - 1$}{
                $\text{set\_bit}(B, j)$ \\
            }
            $\texttt{free\_frames} \leftarrow \texttt{free\_frames} - k$ \\
            \Return $\texttt{run\_start} \times 4096$
        }
    }{
        $\texttt{run} \leftarrow 0$
    }
}
\Return \texttt{None}
\end{algorithm}

La complejidad es $\bigO(N)$ en el peor caso (escaneo lineal completo).

\subsection{Tabla de Refcount para CoW}

Para soportar Copy-on-Write, el módulo \texttt{mm} mantiene una tabla global de
contadores de referencia:

\[
\texttt{FRAME\_REFCOUNT}[i] \in \{0, 1, 2, \ldots, 65535\}, \quad i = \texttt{phys} \gg 12
\]

\begin{itemize}
    \item $\texttt{rc} = 0$: marco sin rastreo (propietario único, caso normal).
    \item $\texttt{rc} = 1$: compartido pero con una sola referencia restante
          (puede hacerse escribible sin copia).
    \item $\texttt{rc} \geq 2$: compartido, se requiere copia CoW al escribir.
\end{itemize}

La tabla soporta hasta 524\,288 marcos ($= 2$\,GiB de RAM), usando 1\,MiB
de memoria estática (\texttt{[u16; 524288]}).

% --------------------------------------------------------------------
\section{Tablas de Páginas: x86-64 de Cuatro Niveles}
\label{sec:page-tables}

\subsection{Estructura de Traducción}

La arquitectura x86-64 utiliza tablas de páginas de cuatro niveles para traducir
direcciones virtuales de 48 bits a direcciones físicas:

\begin{figure}[ht]
\centering
\begin{tikzpicture}[
    level/.style={
        draw, minimum width=3cm, minimum height=0.8cm,
        fill=blue!10, font=\ttfamily\small
    },
    arrow/.style={-{Stealth[length=4pt]}, thick},
    bits/.style={font=\scriptsize\ttfamily, text=red!70!black},
    node distance=0.5cm
]

% Virtual address decomposition
\node[draw, minimum width=14cm, minimum height=0.8cm, fill=yellow!15] (va) at (0, 5)
    {};
\node[bits] at (-5.25, 5) {Sign ext};
\node[bits] at (-3.15, 5) {PML4 [47:39]};
\node[bits] at (-0.7, 5) {PDP [38:30]};
\node[bits] at (1.75, 5) {PD [29:21]};
\node[bits] at (4.0, 5) {PT [20:12]};
\node[bits] at (6.0, 5) {Offset [11:0]};

% Draw separators in virtual address
\draw[thick] (-6.5, 4.6) -- (-6.5, 5.4);
\draw[thick] (-4.0, 4.6) -- (-4.0, 5.4);
\draw[thick] (-1.7, 4.6) -- (-1.7, 5.4);
\draw[thick] (0.6, 4.6) -- (0.6, 5.4);
\draw[thick] (2.9, 4.6) -- (2.9, 5.4);
\draw[thick] (5.0, 4.6) -- (5.0, 5.4);

% Page table levels
\node[level] (pml4) at (-4.5, 3) {PML4};
\node[level] (pdp) at (-1.5, 3) {PDP};
\node[level] (pd) at (1.5, 3) {PD};
\node[level] (pt) at (4.5, 3) {PT};

% Physical frame
\node[draw, minimum width=2.5cm, minimum height=0.8cm, fill=green!15,
      font=\ttfamily\small] (frame) at (4.5, 1) {Marco 4\,KiB};

% CR3 register
\node[draw, minimum width=2cm, minimum height=0.6cm, fill=red!15,
      font=\ttfamily\small] (cr3) at (-6.5, 3) {CR3};

% Arrows
\draw[arrow] (cr3) -- (pml4);
\draw[arrow] (pml4) -- node[above, font=\scriptsize] {phys addr} (pdp);
\draw[arrow] (pdp) -- node[above, font=\scriptsize] {phys addr} (pd);
\draw[arrow] (pd) -- node[above, font=\scriptsize] {phys addr} (pt);
\draw[arrow] (pt) -- node[right, font=\scriptsize] {phys addr} (frame);

% Index arrows from VA
\draw[arrow, dashed, gray] (-3.15, 4.6) -- (pml4.north);
\draw[arrow, dashed, gray] (-0.7, 4.6) -- (pdp.north);
\draw[arrow, dashed, gray] (1.75, 4.6) -- (pd.north);
\draw[arrow, dashed, gray] (4.0, 4.6) -- (pt.north);

% Labels
\node[font=\small, anchor=west] at (-7, 5.8) {Dirección virtual (48 bits):};
\node[font=\scriptsize, text=gray] at (-4.5, 2.3) {512 entradas};
\node[font=\scriptsize, text=gray] at (-1.5, 2.3) {512 entradas};
\node[font=\scriptsize, text=gray] at (1.5, 2.3) {512 entradas};
\node[font=\scriptsize, text=gray] at (4.5, 2.3) {512 entradas};

\end{tikzpicture}
\caption{Traducción de direcciones virtuales x86-64 con tablas de páginas de cuatro niveles.
Cada nivel tiene 512 entradas de 8 bytes (4\,KiB por tabla). Los bits de índice de la
dirección virtual seleccionan la entrada en cada nivel.}
\label{fig:page-walk}
\end{figure}

Cada entrada de tabla de páginas (PTE) tiene 64 bits con la siguiente estructura:

\begin{center}
\begin{tabular}{@{}clp{8cm}@{}}
\toprule
\textbf{Bit(s)} & \textbf{Nombre} & \textbf{Descripción} \\
\midrule
0     & \texttt{PRESENT}       & Entrada válida (si $= 0$, acceso genera \#PF) \\
1     & \texttt{WRITABLE}      & Página escribible (si $= 0$, escritura genera \#PF) \\
2     & \texttt{USER}          & Accesible desde Ring~3 \\
3     & \texttt{WRITE\_THROUGH} & Write-through caching \\
4     & \texttt{CACHE\_DISABLE} & Deshabilitar caché (para MMIO) \\
12--51 & Dirección física      & Bits [51:12] de la dirección física del marco \\
63    & \texttt{NO\_EXECUTE}   & Ejecución prohibida (requiere NXE en IA32\_EFER) \\
\bottomrule
\end{tabular}
\end{center}

\subsection{Extracción de Índices}

Dado un virtual address $v$, los índices para cada nivel se calculan como:
\begin{align}
    \text{pml4\_index}(v) &= (v \gg 39) \,\texttt{\&}\, \texttt{0x1FF} \label{eq:pml4-idx} \\
    \text{pdp\_index}(v)  &= (v \gg 30) \,\texttt{\&}\, \texttt{0x1FF} \label{eq:pdp-idx} \\
    \text{pd\_index}(v)   &= (v \gg 21) \,\texttt{\&}\, \texttt{0x1FF} \label{eq:pd-idx} \\
    \text{pt\_index}(v)   &= (v \gg 12) \,\texttt{\&}\, \texttt{0x1FF} \label{eq:pt-idx}
\end{align}

\subsection{Algoritmos de Tabla de Páginas}

\begin{algorithm}[H]
\caption{\textsc{Map-Page}: Mapear una página virtual a un marco físico}
\label{alg:map-page}
\KwIn{Espacio de direcciones $\mathit{AS}$ (PML4 físico), dirección virtual $v$, dirección física $p$, flags $f$}
$\texttt{hhdm} \leftarrow$ offset del Higher Half Direct Map \\
$\texttt{table\_flags} \leftarrow \texttt{PRESENT} \bitor \texttt{WRITABLE} \bitor \texttt{USER}$ \\[4pt]
\tcp{PML4 → PDP}
$\texttt{pml4} \leftarrow (\mathit{AS}.\texttt{pml4\_phys} + \texttt{hhdm})$ como puntero \\
$\texttt{pdp\_phys} \leftarrow \textsc{Ensure-Table}(\texttt{pml4}[\text{pml4\_index}(v)], \texttt{table\_flags})$ \\[4pt]
\tcp{PDP → PD}
$\texttt{pdp} \leftarrow (\texttt{pdp\_phys} + \texttt{hhdm})$ como puntero \\
$\texttt{pd\_phys} \leftarrow \textsc{Ensure-Table}(\texttt{pdp}[\text{pdp\_index}(v)], \texttt{table\_flags})$ \\[4pt]
\tcp{PD → PT}
$\texttt{pd} \leftarrow (\texttt{pd\_phys} + \texttt{hhdm})$ como puntero \\
$\texttt{pt\_phys} \leftarrow \textsc{Ensure-Table}(\texttt{pd}[\text{pd\_index}(v)], \texttt{table\_flags})$ \\[4pt]
\tcp{Escribir entrada hoja}
$\texttt{pt} \leftarrow (\texttt{pt\_phys} + \texttt{hhdm})$ como puntero \\
$\texttt{pt}[\text{pt\_index}(v)] \leftarrow (p \,\texttt{\&}\, \texttt{!0xFFF}) \bitor f$
\end{algorithm}

\begin{algorithm}[H]
\caption{\textsc{Ensure-Table}: Garantizar existencia de tabla intermedia}
\label{alg:ensure-table}
\KwIn{Referencia a entrada $e$, flags $f$}
\KwOut{Dirección física de la tabla apuntada}
\eIf{$e \,\texttt{\&}\, \texttt{PRESENT} \neq 0$}{
    \Return $e \,\texttt{\&}\, \texttt{0x000F\_FFFF\_FFFF\_F000}$ \tcp*{extraer phys}
}{
    $\texttt{frame} \leftarrow \textsc{Alloc-Table-Frame}()$ \tcp*{4\,KiB, puesto a cero}
    $e \leftarrow \texttt{frame} \bitor f$ \\
    \Return \texttt{frame}
}
\end{algorithm}

\begin{algorithm}[H]
\caption{\textsc{Lookup-Phys}: Traducir dirección virtual a física}
\label{alg:lookup-phys}
\KwIn{Espacio de direcciones $\mathit{AS}$, dirección virtual $v$}
\KwOut{Dirección física, o \texttt{None}}
$\texttt{hhdm} \leftarrow$ offset HHDM \\[4pt]
\tcp{PML4}
$\texttt{pml4e} \leftarrow *(\mathit{AS}.\texttt{pml4\_phys} + \texttt{hhdm})[\text{pml4\_index}(v)]$ \\
\If{$\texttt{pml4e} \,\texttt{\&}\, \texttt{PRESENT} = 0$}{\Return \texttt{None}}
$\texttt{pdp\_phys} \leftarrow \texttt{pml4e} \,\texttt{\&}\, \texttt{0x000F\_FFFF\_FFFF\_F000}$ \\[4pt]
\tcp{PDP}
$\texttt{pdpe} \leftarrow *(\texttt{pdp\_phys} + \texttt{hhdm})[\text{pdp\_index}(v)]$ \\
\If{$\texttt{pdpe} \,\texttt{\&}\, \texttt{PRESENT} = 0$}{\Return \texttt{None}}
$\texttt{pd\_phys} \leftarrow \texttt{pdpe} \,\texttt{\&}\, \texttt{0x000F\_FFFF\_FFFF\_F000}$ \\[4pt]
\tcp{PD}
$\texttt{pde} \leftarrow *(\texttt{pd\_phys} + \texttt{hhdm})[\text{pd\_index}(v)]$ \\
\If{$\texttt{pde} \,\texttt{\&}\, \texttt{PRESENT} = 0$}{\Return \texttt{None}}
$\texttt{pt\_phys} \leftarrow \texttt{pde} \,\texttt{\&}\, \texttt{0x000F\_FFFF\_FFFF\_F000}$ \\[4pt]
\tcp{PT (hoja)}
$\texttt{pte} \leftarrow *(\texttt{pt\_phys} + \texttt{hhdm})[\text{pt\_index}(v)]$ \\
\If{$\texttt{pte} \,\texttt{\&}\, \texttt{PRESENT} = 0$}{\Return \texttt{None}}
\Return $\texttt{pte} \,\texttt{\&}\, \texttt{0x000F\_FFFF\_FFFF\_F000}$
\end{algorithm}

\begin{algorithm}[H]
\caption{\textsc{Unmap-Page}: Desmapear una página virtual}
\label{alg:unmap-page}
\KwIn{Espacio de direcciones $\mathit{AS}$, dirección virtual $v$}
\KwOut{\texttt{Some(phys)} si la página existía, \texttt{None} si no}
Recorrer los cuatro niveles como en \textsc{Lookup-Phys} \\
\If{algún nivel intermedio no está presente}{\Return \texttt{None}}
$\texttt{pte} \leftarrow$ referencia mutable a la entrada hoja \\
\If{$\texttt{pte} \,\texttt{\&}\, \texttt{PRESENT} = 0$}{\Return \texttt{None}}
$\texttt{phys} \leftarrow \texttt{pte} \,\texttt{\&}\, \texttt{0x000F\_FFFF\_FFFF\_F000}$ \\
$\texttt{pte} \leftarrow 0$ \tcp*{borrar entrada}
\Return $\texttt{Some(phys)}$
\end{algorithm}

% --------------------------------------------------------------------
\section{Copy-on-Write (CoW)}
\label{sec:cow}

El mecanismo de Copy-on-Write permite que \texttt{fork()} sea eficiente: en lugar
de copiar toda la memoria del proceso padre, ambos procesos comparten las mismas
páginas físicas marcadas como solo lectura. Cuando cualquiera de los dos escribe
en una página compartida, se genera un fallo de página que el VMM resuelve
copiando la página.

\subsection{Algoritmo de Clonación CoW}

\begin{algorithm}[H]
\caption{\textsc{Clone-CoW}: Clonar un espacio de direcciones con Copy-on-Write}
\label{alg:clone-cow}
\KwIn{Espacio de direcciones padre $P$}
\KwOut{Nuevo espacio de direcciones hijo $C$}
$C.\texttt{pml4\_phys} \leftarrow \textsc{Alloc-Table-Frame}()$ \\[4pt]
\tcp{Copiar mitad del kernel (entradas 256..511)}
\For{$i \leftarrow 256$ \KwTo $511$}{
    $C.\texttt{pml4}[i] \leftarrow P.\texttt{pml4}[i]$
}
\tcp{Clonar mitad de usuario (entradas 0..255)}
\For{$i_4 \leftarrow 0$ \KwTo $255$}{
    \If{$P.\texttt{pml4}[i_4]$ no está presente}{\textbf{continuar}}
    Asignar nueva PDP para $C$ \\
    \For{$i_3 \leftarrow 0$ \KwTo $511$}{
        \If{PDP padre $[i_3]$ no está presente}{\textbf{continuar}}
        Asignar nuevo PD para $C$ \\
        \For{$i_2 \leftarrow 0$ \KwTo $511$}{
            \If{PD padre $[i_2]$ no está presente}{\textbf{continuar}}
            Asignar nueva PT para $C$ \\
            \For{$i_1 \leftarrow 0$ \KwTo $511$}{
                $\texttt{pte} \leftarrow$ PT padre $[i_1]$ \\
                \If{$\texttt{pte}$ no está presente}{\textbf{continuar}}
                $\texttt{phys} \leftarrow \texttt{pte} \,\texttt{\&}\, \text{PHYS\_MASK}$ \\[4pt]
                \tcp{Hijo: PTE sin bit WRITABLE}
                $C.\text{PT}[i_1] \leftarrow \texttt{pte} \,\texttt{\&}\, \lnot\texttt{WRITABLE}$ \\[4pt]
                \tcp{Padre: marcar read-only (CoW simétrico)}
                \If{$\texttt{pte}$ era writable $\land$ no es MMIO $\land$ no es región de servicio}{
                    $P.\text{PT}[i_1] \leftarrow \texttt{pte} \,\texttt{\&}\, \lnot\texttt{WRITABLE}$
                }
                \tcp{Refcount: 0→2 (primer compartimiento), >0→+1}
                \eIf{$\texttt{REFCOUNT}[\texttt{phys} \gg 12] = 0$}{
                    $\texttt{REFCOUNT}[\texttt{phys} \gg 12] \leftarrow 2$
                }{
                    $\texttt{REFCOUNT}[\texttt{phys} \gg 12] \mathrel{+}= 1$
                }
            }
        }
    }
}
\tcp{Flush TLB del padre (PTEs cambiaron a read-only)}
$\textsc{Flush-TLB}()$ \\
\Return $C$
\end{algorithm}

\begin{remark}
El CoW es \emph{simétrico}: ambos, padre e hijo, pierden el permiso de escritura.
Esto garantiza que el \emph{primer} escritor (sea padre o hijo) pague el costo de
la copia, no solo el hijo.
Las páginas UC/MMIO y las regiones de servicio del proceso \texttt{init} se
excluyen para no romper DMA ni el handler de IPC.
\end{remark}

\subsection{Manejo de Fallos CoW en el VMM}

\begin{algorithm}[H]
\caption{\textsc{CoW-Fault-Handler}: Resolución de fallo de escritura CoW}
\label{alg:cow-fault}
\KwIn{Fallo de página con código $c$ y dirección $v$, espacio de direcciones $\mathit{AS}$}
\tcp{Detectar fallo CoW: escritura (bit 1) a página presente (bit 0)}
\If{$(c \,\texttt{\&}\, \texttt{0x03}) \neq \texttt{0x03}$}{
    \textbf{no es CoW, manejar como fallo normal}
}
$\texttt{new\_frame} \leftarrow \textsc{Frame-Alloc}()$ \\
\tcp{Copiar 4\,KiB del marco original al nuevo (vía HHDM en kernel)}
$\textsc{Frame-Copy}(\texttt{new\_frame}, \mathit{AS}, v)$ \\
\tcp{Reemplazar el mapeo: desmapear viejo, mapear nuevo como WRITABLE}
$\textsc{Unmap-From}(\mathit{AS}, v)$ \\
$\textsc{Map-Into}(\mathit{AS}, v, \texttt{new\_frame}, \texttt{WRITABLE})$ \\
$\textsc{Thread-Resume}(\texttt{tid})$
\end{algorithm}

\begin{theorem}[CoW preserva aislamiento de memoria]
\label{thm:cow-isolation}
Después de una operación \textsc{Clone-CoW}, toda escritura del padre o del hijo
a una página compartida resulta en una copia privada. Formalmente: sean
$M_P$ y $M_C$ las vistas de memoria del padre y del hijo respectivamente.
Para toda dirección virtual $v$ en el espacio compartido, después de una escritura
$M_P[v] \leftarrow x$ o $M_C[v] \leftarrow x$, los marcos físicos subyacentes
de $M_P[v]$ y $M_C[v]$ son distintos.
\end{theorem}

\begin{proof}
Sea $v$ una dirección virtual compartida tras \textsc{Clone-CoW}. Sea $f_0$ el
marco físico original.

\textbf{Invariante}: Tras \textsc{Clone-CoW}, la PTE de $v$ en \emph{ambos}
espacios de direcciones tiene el bit \texttt{WRITABLE} desactivado.

\textbf{Caso 1: El padre escribe en $v$.} La escritura genera un \#PF con
código $c$ donde $c \,\texttt{\&}\, \texttt{0x03} = \texttt{0x03}$ (escritura a
página presente). El VMM ejecuta \textsc{CoW-Fault-Handler}:
\begin{enumerate}
    \item Asigna un nuevo marco $f_P$.
    \item Copia el contenido de $f_0$ a $f_P$ (4\,KiB).
    \item Desmapea $v \to f_0$ en el padre y mapea $v \to f_P$ con \texttt{WRITABLE}.
\end{enumerate}
Ahora $M_P[v]$ reside en $f_P$ y $M_C[v]$ sigue en $f_0$. Son marcos distintos:
$f_P \neq f_0$ porque \textsc{Frame-Alloc} nunca retorna un marco ya asignado
(por el \cref{inv:bitmap-consistency}).

\textbf{Caso 2: El hijo escribe en $v$.} Análogo al Caso~1, con los roles
invertidos. El hijo obtiene un nuevo marco $f_C \neq f_0$.

\textbf{Caso 3: Ambos escriben.} Después de la primera escritura (digamos del
padre), tenemos $M_P[v] \to f_P$ y $M_C[v] \to f_0$. Cuando el hijo escribe,
si $f_0$ aún tiene \texttt{REFCOUNT} $\geq 2$, se genera otro fallo CoW
y el hijo obtiene $f_C \neq f_0$. En total, $f_P$, $f_C$, y $f_0$ (si aún
existe) son marcos distintos.

En todos los casos, la escritura resulta en marcos físicos independientes, lo
que preserva el aislamiento. \qedhere
\end{proof}

\begin{theorem}[Cota de memoria para CoW]
\label{thm:cow-memory}
Sea $n$ el número de páginas del proceso original. Después de un \texttt{fork()}
con CoW:
\begin{enumerate}
    \item La memoria total inmediatamente después del fork es $n$ páginas
          (compartidas) más $\bigO(n/512)$ páginas para las tablas de páginas
          del hijo (una PT por cada 512 entradas presentes).
    \item En el peor caso, si ambos procesos escriben en todas las páginas,
          la memoria total es $\leq 2n$ páginas de datos más tablas.
\end{enumerate}
\end{theorem}

\begin{proof}
\begin{enumerate}
    \item Inmediatamente tras \textsc{Clone-CoW}, no se duplica ninguna página
          de datos: los marcos físicos son compartidos. Solo se asignan nuevas
          tablas de páginas (PML4, PDP, PD, PT) para el hijo. En el peor caso,
          si las $n$ páginas están dispersas en $n$ PTs distintas, se necesitan
          $n$ PTs + hasta $n/512$ PDs + hasta $n/512^2$ PDPs + 1 PML4.
          Esto es $\bigO(n/512)$ en el caso típico donde las páginas están
          agrupadas.

    \item Cada escritura CoW duplica exactamente una página (4\,KiB). Si el padre
          escribe en las $n$ páginas, genera $n$ copias. El hijo conserva las $n$
          originales. Total: $2n$ páginas de datos. Análogamente si solo el hijo
          escribe, o si ambos escriben (máximo $2n$ en cualquier caso). \qedhere
\end{enumerate}
\end{proof}

% --------------------------------------------------------------------
\section{Enforcement W\^{}X via mprotect}
\label{sec:wx}

\sotos{} aplica la política Write XOR Execute (W\^{}X) a nivel de PTE: una página
no puede tener simultáneamente los bits \texttt{WRITABLE} y la ausencia de
\texttt{NO\_EXECUTE}. La syscall \texttt{SYS\_PROTECT} (134) permite cambiar los
permisos de una página existente con esta restricción.

\begin{definition}[Propiedad W\^{}X]
Para toda PTE $e$ con $e \,\texttt{\&}\, \texttt{PRESENT} \neq 0$ en un espacio
de direcciones sin relajación:
\[
(e \,\texttt{\&}\, \texttt{WRITABLE} \neq 0) \implies (e \,\texttt{\&}\, \texttt{NO\_EXECUTE} \neq 0)
\]
Equivalentemente, si una página es escribible, no es ejecutable, y viceversa.
\end{definition}

Para binarios que requieren páginas RWX (como Wine/glibc con IAT fixups o JIT),
se ofrece un mecanismo de relajación por espacio de direcciones mediante un bitmap
atómico \texttt{WX\_RELAXED}. Cuando está activado, el kernel no añade
\texttt{NO\_EXECUTE} automáticamente a páginas escribibles de ese AS.

% ====================================================================
\chapter{LUCAS --- Capa de Compatibilidad Linux}
\label{chap:lucas}
% ====================================================================

LUCAS (\emph{Linux User Compatibility Abstraction Shim}) es la capa de
espacio de usuario que permite ejecutar binarios de Linux sin modificación sobre
\sotos{}. Intercepta las llamadas al sistema Linux del proceso huésped, las traduce
a operaciones sobre los servicios del micronúcleo, y devuelve resultados en el
formato esperado por la ABI de Linux.

% --------------------------------------------------------------------
\section{Arquitectura del Redireccionamiento de Syscalls}
\label{sec:lucas-arch}

Cuando un proceso huésped ejecuta la instrucción \texttt{syscall}, el kernel de
\sotos{} detecta que el hilo tiene un \emph{endpoint de redirección} configurado
y, en lugar de manejar la syscall directamente, la entrega como un mensaje IPC
al manejador LUCAS.

\begin{figure}[ht]
\centering
\begin{tikzpicture}[
    box/.style={draw, rounded corners=3pt, minimum width=2.8cm,
                minimum height=1cm, font=\small, align=center},
    arrow/.style={-{Stealth[length=4pt]}, thick},
    dasharrow/.style={-{Stealth[length=4pt]}, thick, dashed},
    node distance=1.5cm
]

% Nodes
\node[box, fill=orange!15] (guest) {Proceso\\Huésped};
\node[box, fill=red!15, right=2cm of guest] (kernel) {Kernel\\(\texttt{syscall})};
\node[box, fill=blue!15, right=2cm of kernel] (lucas) {LUCAS\\Handler};
\node[box, fill=green!15, below=1.5cm of lucas] (service) {Servicio\\(Net/VFS/VMM)};

% Forward path
\draw[arrow] (guest) -- node[above, font=\scriptsize] {\texttt{syscall}} (kernel);
\draw[arrow] (kernel) -- node[above, font=\scriptsize] {IPC msg} node[below, font=\scriptsize] {(regs $\to$ msg)} (lucas);
\draw[dasharrow] (lucas) -- node[right, font=\scriptsize] {IPC a servicio} (service);

% Return path
\draw[arrow, blue!60] (service.west) -| node[below left, font=\scriptsize, pos=0.3] {reply} ([xshift=-5mm]lucas.south);
\draw[arrow, blue!60] (lucas.north west) -- ++(0, 0.5) -| node[above, font=\scriptsize, pos=0.7] {IPC reply} (kernel.north);
\draw[arrow, blue!60] (kernel.south) -- ++(0, -0.5) -| node[below, font=\scriptsize, pos=0.3] {\texttt{sysret}} (guest.south);

\end{tikzpicture}
\caption{Flujo de una syscall Linux en \sotos{}. El kernel redirige la syscall como
mensaje IPC a LUCAS, que opcionalmente delega a un servicio del sistema, y devuelve
el resultado al huésped vía IPC reply y \texttt{SYSRET}.}
\label{fig:lucas-flow}
\end{figure}

\subsection{Recorrido Completo de una Syscall: \texttt{write(1, "Hola", 4)}}

Para ilustrar el mecanismo, rastreamos una llamada \texttt{write(fd=1, buf="Hola", len=4)}
paso a paso:

\begin{enumerate}
    \item \textbf{Huésped}: El binario Linux coloca \texttt{RAX=1} (SYS\_write),
          \texttt{RDI=1} (fd), \texttt{RSI=ptr} (buffer), \texttt{RDX=4} (longitud)
          y ejecuta \texttt{syscall}.

    \item \textbf{Kernel (entrada)}: El handler de \texttt{SYSCALL} detecta que el
          hilo tiene \texttt{redirect\_ep $\neq$ 0}. Guarda los registros del huésped
          en un \texttt{TrapFrame} y construye un mensaje IPC con los valores de
          RAX, RDI, RSI, RDX, R10, R8, R9. Entrega el mensaje al endpoint de
          redirección (IPC call bloqueante).

    \item \textbf{LUCAS (child\_handler)}: Recibe el mensaje IPC. Extrae
          \texttt{syscall\_nr = regs[0]} ($= 1$, SYS\_WRITE). El manejador del módulo
          \texttt{fs.rs} lee 4 bytes desde la dirección \texttt{RSI} en el espacio de
          memoria del huésped (usando \texttt{SYS\_GET\_THREAD\_REGS} para acceso remoto
          o memoria compartida). Si \texttt{fd=1} (stdout), escribe los bytes a la
          consola serial vía \texttt{sys::debug\_print()}.

    \item \textbf{LUCAS (reply)}: Construye el mensaje de respuesta con
          \texttt{regs[0] = 4} (bytes escritos), y usa la etiqueta
          \texttt{SIG\_REDIRECT\_TAG} para restaurar el RIP/RSP/RAX del huésped.
          Responde por IPC.

    \item \textbf{Kernel (retorno)}: Recibe el reply, aplica el
          \texttt{SIG\_REDIRECT\_TAG}: establece el RIP, RSP y RAX del huésped
          según los registros del reply. Ejecuta \texttt{SYSRET} para devolver
          el control al huésped.

    \item \textbf{Huésped}: Reanuda ejecución después del \texttt{syscall} con
          \texttt{RAX=4} (éxito, 4 bytes escritos).
\end{enumerate}

% --------------------------------------------------------------------
\section{Emulación de la ABI Linux: 127+ Syscalls}
\label{sec:lucas-abi}

El manejador LUCAS implementa más de 127 llamadas al sistema Linux, organizadas
en seis módulos:

\begin{table}[ht]
\centering
\caption{Módulos del manejador LUCAS y syscalls representativas.}
\label{tab:lucas-modules}
\begin{tabular}{@{}llp{7cm}@{}}
\toprule
\textbf{Módulo} & \textbf{Archivo} & \textbf{Syscalls principales} \\
\midrule
Sistema de archivos & \texttt{fs.rs} & open, openat, read, write, writev, close,
    lseek, fstat, fstatat, getdents64, mkdir, unlink, rename, dup/dup2/dup3, pipe,
    fcntl, ioctl, mmap (file-backed), pread64, pwrite64, preadv, pwritev \\
Memoria & \texttt{mm.rs} & brk, mmap, munmap, mprotect, mremap, madvise, memfd\_create \\
Red & \texttt{net.rs} & socket, connect, bind, sendto, recvfrom, recvmsg, readv,
    poll, epoll\_create/ctl/wait, select, pselect6 \\
Señales & \texttt{signal.rs} & rt\_sigaction, rt\_sigprocmask, rt\_sigreturn,
    rt\_sigpending, rt\_sigtimedwait, rt\_sigsuspend, sigaltstack, tgkill \\
Procesos & \texttt{task.rs} & fork, clone, execve, wait4, waitid, exit, exit\_group,
    getpid, gettid, set\_tid\_address, sched\_yield \\
Información & \texttt{info.rs} & uname, getuid/gid/euid/egid, clock\_gettime,
    clock\_getres, gettimeofday, getrandom, getrusage, prlimit64, sysinfo \\
\bottomrule
\end{tabular}
\end{table}

% --------------------------------------------------------------------
\section{Modelo de Procesos: Fork/Exec/Wait}
\label{sec:lucas-process}

LUCAS emula el modelo de procesos POSIX completo:

\subsection{Grupos de Hilos (\texttt{CLONE\_THREAD})}

Cada proceso tiene un \emph{thread group ID} (TGID) y puede crear hilos
adicionales mediante \texttt{clone()} con el flag \texttt{CLONE\_THREAD}.
Los hilos de un grupo comparten:

\begin{itemize}
    \item \textbf{Tabla de descriptores de archivo}: via \texttt{CLONE\_FILES}.
          Las tablas están en arreglos estáticos por grupo (\texttt{GRP\_FDS},
          \texttt{GRP\_INITRD}, \texttt{GRP\_VFS}).
    \item \textbf{Espacio de direcciones}: via \texttt{CLONE\_VM}. El hilo hijo
          comparte el mismo CR3.
    \item \textbf{Manejadores de señal}: via \texttt{CLONE\_SIGHAND}.
\end{itemize}

\subsection{Fork con CoW}

La llamada \texttt{fork()} invoca \texttt{SYS\_AS\_CLONE} (125) en el kernel,
que ejecuta \textsc{Clone-CoW} (\cref{alg:clone-cow}). El hijo se configura
con un trampolín en el vDSO que establece \texttt{RAX=0} (retorno de fork) y
restaura los registros callee-saved:

\begin{lstlisting}[language={[x86masm]Assembler},caption={Trampolín CoW fork child restore (vDSO 0xB80740).},label={lst:cow-trampoline}]
cow_fork_restore:
    xor eax, eax        ; fork() retorna 0 en el hijo
    pop rbx
    pop rbp
    pop r12
    pop r13
    pop r14
    pop r15
    ret                  ; retorna al punto de fork del padre
\end{lstlisting}

\subsection{Entrega de Señales}

Las señales se entregan en dos modos:

\begin{enumerate}
    \item \textbf{Síncrono} (en frontera de syscall): Si hay una señal pendiente
          cuando LUCAS devuelve un reply, modifica el RIP del huésped para saltar
          al manejador de señal antes de reanudar la ejecución normal.

    \item \textbf{Asíncrono} (vía interrupción de timer): El kernel detecta señales
          pendientes en la rutina de interrupción del LAPIC timer (cada $\sim$10\,ms).
          Si el hilo tiene \texttt{pending\_signals $\neq$ 0} y un
          \texttt{signal\_trampoline} registrado, el kernel:
          \begin{enumerate}
              \item Guarda el \texttt{TrapFrame} actual del huésped en un
                    \texttt{SignalFrame} de 176 bytes en la pila del huésped.
              \item Redirige RIP al trampolín de señales en el vDSO (0xB80700).
              \item El trampolín invoca el handler registrado por \texttt{sigaction()}.
              \item Al retornar, el handler ejecuta \texttt{rt\_sigreturn} (syscall 15),
                    que restaura el \texttt{TrapFrame} original desde el
                    \texttt{SignalFrame}.
          \end{enumerate}
\end{enumerate}

% --------------------------------------------------------------------
\section{Puente al Sistema de Archivos Virtual}
\label{sec:lucas-vfs}

\subsection{ObjectStore: Almacenamiento Transaccional}

El sistema de archivos de \sotos{} se basa en el \texttt{ObjectStore}, un almacén
de objetos sobre disco con escritura adelantada (WAL) para consistencia ante
fallos. Los detalles del layout en disco se describen en el \cref{sec:objstore}.

\subsection{Resolución de Rutas con CWD}

Cada proceso mantiene un directorio de trabajo actual (\texttt{CWD}) en el arreglo
\texttt{GRP\_CWD}. La función \texttt{resolve\_with\_cwd()} resuelve rutas relativas
concatenándolas con el CWD:

\begin{itemize}
    \item Si la ruta comienza con \texttt{/}, se usa tal cual (ruta absoluta).
    \item Si no, se concatena \texttt{CWD + "/" + ruta}.
    \item Se aplica a todas las syscalls de archivo: open, openat, stat, access,
          mkdir, unlink, rename.
\end{itemize}

\subsection{Snapshots CoW}

El \texttt{ObjectStore} soporta hasta 4 snapshots simultáneos. Cada snapshot
captura el estado del directorio y bitmap en un momento dado, permitiendo
lecturas consistentes mientras las escrituras continúan en la copia activa.

% --------------------------------------------------------------------
\section{vDSO: Shared Library Forjado}
\label{sec:vdso}

El \emph{Virtual Dynamic Shared Object} (vDSO) es una página de 4\,KiB mapeada
en el espacio de direcciones de cada proceso en \texttt{0xB80000}. Contiene un
ELF mínimo pero completo (\texttt{ET\_DYN}) que exporta funciones de tiempo
accesibles sin transición a Ring~0.

\begin{figure}[ht]
\centering
\begin{tikzpicture}[
    block/.style={draw, minimum width=6cm, minimum height=0.6cm,
                  font=\ttfamily\scriptsize, anchor=north west},
    label/.style={font=\scriptsize\ttfamily, anchor=east},
    node distance=0
]

% Page boundary
\draw[thick] (0, 0) rectangle (6.5, -8);
\node[font=\small\bfseries] at (3.25, 0.3) {Página vDSO (4\,KiB @ 0xB80000)};

% Sections
\node[block, fill=red!10] (ehdr) at (0.25, -0.1) {ELF Header (64\,B)};
\node[block, fill=red!15, below=0pt of ehdr] (phdr) {Program Headers (2 $\times$ PHDR)};
\node[label] at (-0.1, -0.4) {0x040};

\node[block, fill=yellow!15, below=3pt of phdr] (dyn) {.dynamic};
\node[label] at (-0.1, -1.2) {0x100};

\node[block, fill=blue!10, below=3pt of dyn] (dynsym) {.dynsym (5 $\times$ Elf64\_Sym)};
\node[label] at (-0.1, -1.9) {0x200};

\node[block, fill=blue!15, below=3pt of dynsym] (dynstr) {.dynstr};
\node[label] at (-0.1, -2.6) {0x280};

\node[block, fill=purple!10, below=3pt of dynstr] (hash) {.hash (SysV)};
\node[label] at (-0.1, -3.3) {0x300};

\node[block, fill=purple!15, below=3pt of hash] (versym) {.gnu.version};
\node[label] at (-0.1, -4.0) {0x340};

\node[block, fill=purple!20, below=3pt of versym] (verdef) {.gnu.version\_d (LINUX\_2.6)};
\node[label] at (-0.1, -4.7) {0x350};

\node[block, fill=green!15, below=3pt of verdef, minimum height=1.8cm] (text) {%
    .text: \_\_vdso\_clock\_gettime\\
    signal\_trampoline (0x700)\\
    sigreturn\_restorer (0x710)\\
    fork\_restore (0x720)\\
    cow\_fork\_restore (0x740)};
\node[label] at (-0.1, -5.4) {0x600};

\node[block, fill=gray!15, below=3pt of text] (data) {.data: boot\_tsc};
\node[label] at (-0.1, -7.65) {0xE00};

\end{tikzpicture}
\caption{Layout de la página vDSO. El ELF es autocontenido en una única página de 4\,KiB.
Las secciones \texttt{.gnu.version} y \texttt{.gnu.version\_d} con la cadena
\texttt{"LINUX\_2.6"} satisfacen la verificación de versión de musl y glibc.}
\label{fig:vdso-layout}
\end{figure}

La función principal exportada, \texttt{\_\_vdso\_clock\_gettime}, usa \texttt{RDTSC}
para devolver el tiempo sin transición al kernel:

\begin{lstlisting}[language={[x86masm]Assembler},caption={Implementación de \texttt{\_\_vdso\_clock\_gettime} en el vDSO.},label={lst:vdso-clock}]
__vdso_clock_gettime:
    rdtsc                    ; EDX:EAX = timestamp counter
    shl rdx, 32
    or  rax, rdx             ; RAX = TSC completo (64 bits)
    ; Dividir por 2 GHz para obtener segundos
    mov rcx, 2000000000
    xor edx, edx
    div rcx                  ; RAX = segundos, RDX = nanosegundos residuales
    mov [rsi],    rax        ; tv_sec
    ; Convertir residuo a nanosegundos
    imul rdx, rdx, 1000000000
    mov rcx2, 2000000000
    div rcx                  ; RDX = nanosegundos
    mov [rsi+8],  rax        ; tv_nsec
    xor eax, eax             ; retornar 0 (exito)
    ret
\end{lstlisting}

El vDSO también contiene trampolines críticos para el funcionamiento del sistema:

\begin{center}
\begin{tabular}{@{}lll@{}}
\toprule
\textbf{Offset} & \textbf{Dirección} & \textbf{Función} \\
\midrule
\texttt{0x700} & \texttt{0xB80700} & Trampolín de señales \\
\texttt{0x710} & \texttt{0xB80710} & Restaurador \texttt{rt\_sigreturn} \\
\texttt{0x720} & \texttt{0xB80720} & Fork restore (pop callee-saved + ret) \\
\texttt{0x730} & \texttt{0xB80730} & Exit stub (\texttt{exit\_group}) \\
\texttt{0x740} & \texttt{0xB80740} & CoW fork restore (xor eax + pop + ret) \\
\texttt{0x760} & \texttt{0xB80760} & Fork TLS trampoline (arch\_prctl + jump) \\
\texttt{0x7C0} & \texttt{0xB807C0} & Pre-TLS entry trampoline \\
\texttt{0x7E0} & \texttt{0xB807E0} & Yield loop (mantiene hilo vivo) \\
\bottomrule
\end{tabular}
\end{center}

% ====================================================================
\chapter{Servicios del Sistema}
\label{chap:servicios}
% ====================================================================

En la arquitectura de \sotos{}, los servicios del sistema ejecutan en espacio
de usuario como procesos independientes, cada uno con su propio espacio de
direcciones (CR3). Se comunican con el kernel y entre sí exclusivamente
mediante IPC (descrito en el \cref{chap:ipc}).

% --------------------------------------------------------------------
\section{Servicio de Red (smoltcp)}
\label{sec:net-service}

El servicio de red implementa la pila TCP/IP completa usando la biblioteca
\texttt{smoltcp}, proporcionando conectividad de red a los procesos huésped
a través de una interfaz IPC de 13 comandos.

\subsection{Interfaz de Comandos IPC}

\begin{table}[ht]
\centering
\caption{Comandos IPC del servicio de red.}
\label{tab:net-commands}
\begin{tabular}{@{}clp{7.5cm}@{}}
\toprule
\textbf{Código} & \textbf{Comando} & \textbf{Descripción} \\
\midrule
1  & \texttt{CMD\_PING}          & Enviar ICMP echo request y esperar reply \\
2  & \texttt{CMD\_DNS\_QUERY}    & Resolver nombre de dominio vía UDP (10.0.2.3) \\
3  & \texttt{CMD\_TCP\_CONNECT}  & Establecer conexión TCP (3-way handshake) \\
4  & \texttt{CMD\_TCP\_SEND}     & Enviar datos por conexión TCP \\
5  & \texttt{CMD\_TCP\_RECV}     & Recibir datos de conexión TCP \\
6  & \texttt{CMD\_TCP\_CLOSE}    & Cerrar conexión TCP (FIN) \\
7  & \texttt{CMD\_TRACEROUTE\_HOP} & Enviar paquete con TTL específico \\
8  & \texttt{CMD\_UDP\_BIND}     & Vincular puerto UDP efímero \\
9  & \texttt{CMD\_UDP\_SENDTO}   & Enviar datagrama UDP \\
10 & \texttt{CMD\_UDP\_RECV}     & Recibir datagrama UDP \\
11 & \texttt{CMD\_TCP\_STATUS}   & Consultar estado de conexión TCP \\
12 & \texttt{CMD\_NET\_MIRROR}   & Activar mirror de paquetes (debug) \\
13 & \texttt{CMD\_UDP\_HAS\_DATA} & Verificar si hay datos UDP pendientes \\
\bottomrule
\end{tabular}
\end{table}

\subsection{Flujo de una Transferencia TCP}

\begin{figure}[ht]
\centering
\begin{tikzpicture}[
    box/.style={draw, rounded corners=3pt, minimum width=2.2cm,
                minimum height=0.8cm, font=\scriptsize, align=center},
    arrow/.style={-{Stealth[length=3pt]}, thick},
    node distance=1.2cm
]

\node[box, fill=orange!15] (guest) {Proceso\\Huésped};
\node[box, fill=blue!15, right=1cm of guest] (proxy) {Socket\\Proxy\\(LUCAS)};
\node[box, fill=red!15, right=1cm of proxy] (ipc) {IPC\\Channel};
\node[box, fill=green!15, right=1cm of ipc] (net) {Net\\Service\\(smoltcp)};
\node[box, fill=purple!15, right=1cm of net] (virtio) {virtio-net\\Driver};
\node[box, fill=gray!15, right=1cm of virtio] (slirp) {QEMU\\SLIRP};

\draw[arrow] (guest) -- node[above, font=\tiny] {syscall} (proxy);
\draw[arrow] (proxy) -- node[above, font=\tiny] {msg} (ipc);
\draw[arrow] (ipc) -- node[above, font=\tiny] {cmd} (net);
\draw[arrow] (net) -- node[above, font=\tiny] {pkt} (virtio);
\draw[arrow] (virtio) -- node[above, font=\tiny] {frame} (slirp);

% Return path
\draw[arrow, blue!50] (slirp.south) -- ++(0, -0.5) -| (guest.south);
\node[font=\tiny, blue!50] at (3.5, -1.2) {respuesta (ruta inversa)};

\end{tikzpicture}
\caption{Ruta completa de una operación de red: desde la syscall del huésped hasta
el backend SLIRP de QEMU, pasando por el proxy de sockets LUCAS, el canal IPC,
el servicio de red con smoltcp, y el driver virtio-net.}
\label{fig:net-path}
\end{figure}

\subsection{Transferencia TCP en Bloques}

El transporte IPC de \sotos{} tiene un límite inherente de datos por mensaje:
40 bytes para envío (\texttt{regs[3..]} de 5 registros $\times$ 8 bytes) y 64 bytes
para recepción (\texttt{regs[0..7]} de 8 registros $\times$ 8 bytes). Para
transferencias TCP grandes (e.g., descargar una página web con \texttt{wget}),
LUCAS fragmenta los datos en chunks de 64 bytes y los reensambla:

\begin{algorithm}[H]
\caption{\textsc{TCP-Recv-Multi}: Recepción TCP multi-chunk}
\label{alg:tcp-recv}
\KwIn{Conexión $\texttt{conn\_id}$, buffer destino $\texttt{buf}$, tamaño $\texttt{len}$}
\KwOut{Bytes recibidos}
$\texttt{total} \leftarrow 0$ \\
\While{$\texttt{total} < \texttt{len}$}{
    $\texttt{msg} \leftarrow \textsc{IPC-Call}(\texttt{CMD\_TCP\_RECV}, \texttt{conn\_id}, \texttt{total})$ \\
    $\texttt{chunk\_len} \leftarrow \texttt{msg.tag}$ \\
    \If{$\texttt{chunk\_len} = 0$}{\textbf{salir} \tcp*{EOF o no más datos}}
    Copiar $\min(\texttt{chunk\_len}, 64)$ bytes de $\texttt{msg.regs[0..7]}$ a $\texttt{buf}[\texttt{total}..]$ \\
    $\texttt{total} \mathrel{+}= \texttt{chunk\_len}$
}
\Return \texttt{total}
\end{algorithm}

\subsection{Resolución DNS por UDP}

El servicio de red resuelve nombres de dominio enviando consultas DNS por UDP
al servidor configurado por DHCP (por defecto \texttt{10.0.2.3}, el DNS de SLIRP
de QEMU). Las respuestas DNS mayores a 64 bytes se recuperan en múltiples chunks
con un offset en el argumento 2.

% --------------------------------------------------------------------
\section{Servicio VMM (Virtual Memory Manager)}
\label{sec:vmm-service}

El servidor VMM es el primer servicio lanzado y maneja todos los fallos de
página del sistema. Su diseño es intencionalmente simple: un bucle infinito
que espera notificaciones de fallo y las resuelve.

\subsection{Manejo de Fallos de Página}

El VMM se registra de dos maneras:
\begin{enumerate}
    \item \textbf{Per-AS}: \texttt{fault\_register\_as(notify, init\_as\_cap)} para
          los fallos del proceso init.
    \item \textbf{Global}: \texttt{fault\_register(notify)} con \texttt{rsi=0},
          que registra al VMM como manejador de respaldo para \emph{cualquier}
          espacio de direcciones (incluyendo los clones CoW de procesos hijos).
\end{enumerate}

Los tipos de fallo que el VMM maneja son:

\begin{enumerate}
    \item \textbf{Demanda de paginación}: Página no presente $\to$ asignar marco,
          mapear como WRITABLE, reanudar hilo.
    \item \textbf{CoW}: Escritura a página presente sin permiso de escritura
          ($\texttt{code} \,\texttt{\&}\, \texttt{0x03} = \texttt{0x03}$) $\to$
          ejecutar \textsc{CoW-Fault-Handler} (\cref{alg:cow-fault}).
    \item \textbf{Violación NX}: Fetch de instrucción en página NX
          ($\texttt{code} \,\texttt{\&}\, \texttt{0x11} = \texttt{0x11}$) $\to$
          cambiar permisos a R+X (sin WRITABLE para mantener W\^{}X).
    \item \textbf{Bucle infinito}: Si la misma dirección genera $> 8$ fallos
          consecutivos $\to$ inyectar SIGSEGV (señal 11) al hilo.
\end{enumerate}

\subsection{Registro Global de Fallos}

\begin{invariant}[Cobertura de fallos]
Todo fallo de página en espacio de usuario es manejado por el VMM. Si un
espacio de direcciones tiene un manejador per-AS registrado, se usa ese.
De lo contrario, el manejador global (CR3 $= 0$) captura el fallo. No
existen fallos de página de usuario sin manejador.
\end{invariant}

% --------------------------------------------------------------------
\section{Servicio de Teclado}
\label{sec:kbd-service}

El driver de teclado PS/2 ejecuta como un proceso independiente que:
\begin{enumerate}
    \item Registra un manejador de IRQ para IRQ~1 (teclado PS/2).
    \item Lee los scancodes del puerto \texttt{0x60}.
    \item Traduce scancodes a eventos de tecla (key down/up).
    \item Escribe los caracteres resultantes en un ring buffer compartido
          en \texttt{0x510000} (dirección centralizada en \texttt{sotos\_common}).
\end{enumerate}

LUCAS lee del ring buffer cuando un proceso huésped ejecuta \texttt{read(0, ...)}
(lectura de stdin, fd~0).

% --------------------------------------------------------------------
\section{ObjectStore y VFS}
\label{sec:objstore}

\subsection{Layout en Disco}

El \texttt{ObjectStore} (versión 6) organiza el disco en regiones contiguas:

\begin{table}[ht]
\centering
\caption{Layout de sectores del ObjectStore v6. Cada sector tiene 512 bytes.}
\label{tab:objstore-layout}
\begin{tabular}{@{}rlll@{}}
\toprule
\textbf{Sector(es)} & \textbf{Región} & \textbf{Tamaño} & \textbf{Contenido} \\
\midrule
0          & Superblock        & 1 sector (512\,B)   & Magia, versión, contadores \\
1          & WAL Header        & 1 sector (512\,B)   & Secuencia, entradas, committed \\
2--5       & WAL Payload       & 4 sectores (2\,KiB) & Copias de sectores a escribir \\
6--133     & Bitmap            & 128 sectores (64\,KiB) & 524\,288 bits = estado de cada sector de datos \\
134--2181  & Directorio        & 2048 sectores       & 8192 entradas de 128\,B \\
2182--2197 & Refcount          & 16 sectores         & 8192 $\times$ \texttt{u8} (para CoW snapshots) \\
2198--2201 & Snapshot Meta     & 4 sectores          & Metadatos de hasta 4 snapshots \\
2202--10393 & Snapshot Dir     & $4 \times 2048$ sectores & Copia del directorio por snapshot \\
10394--10905 & Snapshot Bitmap & $4 \times 128$ sectores & Copia del bitmap por snapshot \\
10906+     & Datos             & Variable            & Contenido de los archivos \\
\bottomrule
\end{tabular}
\end{table}

\subsection{Entrada de Directorio}

Cada entrada de directorio ocupa 128 bytes (4 por sector de 512\,B):

\begin{center}
\begin{tabular}{@{}rll@{}}
\toprule
\textbf{Offset} & \textbf{Campo} & \textbf{Tipo/Tamaño} \\
\midrule
0   & \texttt{oid}            & \texttt{u64} (8\,B) — identificador único del objeto \\
8   & \texttt{name}           & \texttt{[u8; 48]} — nombre, terminado en null \\
56  & \texttt{size}           & \texttt{u64} (8\,B) — tamaño del archivo en bytes \\
64  & \texttt{sector\_start}  & \texttt{u32} (4\,B) — primer sector de datos \\
68  & \texttt{sector\_count}  & \texttt{u32} (4\,B) — número de sectores ocupados \\
72  & \texttt{flags}          & \texttt{u32} (4\,B) — bit 0 = directorio \\
76  & \texttt{created\_time}  & \texttt{u64} (8\,B) — timestamp Unix (segundos) \\
84  & \texttt{modified\_time} & \texttt{u64} (8\,B) — timestamp Unix \\
92  & \texttt{accessed\_time} & \texttt{u64} (8\,B) — timestamp Unix \\
100 & \texttt{parent\_oid}    & \texttt{u64} (8\,B) — OID del directorio padre \\
108 & \texttt{permissions}    & \texttt{u32} (4\,B) — permisos Unix (rwxrwxrwx) \\
112 & \texttt{owner\_uid}     & \texttt{u16} (2\,B) — UID del propietario \\
114 & \texttt{owner\_gid}     & \texttt{u16} (2\,B) — GID del propietario \\
116 & \texttt{\_pad}          & \texttt{[u8; 8]} — relleno a 128 bytes \\
\bottomrule
\end{tabular}
\end{center}

\textbf{Nota sobre alineación}: El campo \texttt{oid} (offset 0) + \texttt{name} (48\,B)
+ \texttt{size} (8\,B) suman 64 bytes, alineando \texttt{sector\_start} en una
frontera de 64\,B. La estructura completa es de 128 bytes exactos, verificado
en tiempo de compilación con \texttt{const\_assert}.

\subsection{Write-Ahead Logging}

El WAL (Write-Ahead Log) garantiza atomicidad de las escrituras. Cada
transacción puede afectar hasta 4 sectores.

\begin{algorithm}[H]
\caption{\textsc{WAL-Write}: Escritura atómica con WAL}
\label{alg:wal-write}
\KwIn{Conjunto de pares $\{(\texttt{target}_i, \texttt{data}_i)\}_{i=1}^{k}$ con $k \leq 4$}
\tcp{Fase 1: Escribir payload al WAL}
\For{$i \leftarrow 1$ \KwTo $k$}{
    Escribir $\texttt{data}_i$ en sector $\texttt{SECTOR\_WAL\_PAYLOAD} + i - 1$
}
\tcp{Fase 2: Escribir header con committed=0}
$\texttt{hdr.seq} \mathrel{+}= 1$ \\
$\texttt{hdr.entry\_count} \leftarrow k$ \\
$\texttt{hdr.targets}[0..k-1] \leftarrow$ sectores destino \\
$\texttt{hdr.committed} \leftarrow 0$ \\
Escribir \texttt{hdr} en sector 1 \\[4pt]
\tcp{Fase 3: Escribir datos a sus destinos finales}
\For{$i \leftarrow 1$ \KwTo $k$}{
    Escribir $\texttt{data}_i$ en sector $\texttt{target}_i$
}
\tcp{Fase 4: Marcar como committed}
$\texttt{hdr.committed} \leftarrow 1$ \\
Escribir \texttt{hdr} en sector 1
\end{algorithm}

\begin{theorem}[Atomicidad del WAL]
\label{thm:wal-atomicity}
Sea $T = \{(\texttt{target}_i, \texttt{data}_i)\}_{i=1}^{k}$ una transacción
ejecutada por \textsc{WAL-Write}. En caso de fallo (corte de energía, reset)
en cualquier punto de la ejecución:
\begin{enumerate}
    \item Si el fallo ocurre antes de la Fase~4 (\texttt{committed} aún $= 0$):
          los sectores destino pueden estar parcialmente escritos, pero el WAL
          aún contiene los datos correctos en el payload. En la recuperación,
          se puede repetir la Fase~3 completa.
    \item Si el fallo ocurre después de la Fase~4: la transacción está
          completamente persistida.
\end{enumerate}
En ambos casos, el sistema de archivos recupera un estado consistente.
\end{theorem}

\begin{proof}
Consideramos los cuatro puntos de fallo posibles:

\textbf{Fallo durante Fase~1} (escritura del payload al WAL): El header aún no
se ha escrito, por lo que en recuperación $\texttt{hdr.entry\_count} = 0$ o el
$\texttt{seq}$ no ha cambiado. No hay transacción pendiente. Los sectores
destino no fueron modificados. Estado: consistente (transacción no ocurrió).

\textbf{Fallo durante Fase~2} (escritura del header): El header puede estar
corrupto o tener \texttt{committed} $= 0$ con \texttt{entry\_count} $> 0$
pero magia incorrecta ($\neq$ \texttt{0x57414C48}). En recuperación, se
verifica la magia; si es inválida, se ignora. Si es válida pero
$\texttt{committed} = 0$, se re-ejecuta Fase~3 con los datos del payload
(que están completos desde Fase~1). Estado: consistente.

\textbf{Fallo durante Fase~3} (escritura a destinos): Algunos sectores
destino pueden estar escritos y otros no. Pero $\texttt{committed} = 0$,
y el payload completo está en el WAL. En recuperación, se re-ejecuta Fase~3
completa (idempotente: escribir el mismo dato al mismo sector). Estado: consistente.

\textbf{Fallo durante Fase~4} (marcar committed): Si la escritura del header
con $\texttt{committed} = 1$ no se completó, en recuperación
$\texttt{committed}$ será 0, y se repite Fase~3 (los datos ya están en los
destinos, pero la repetición es inofensiva). Si se completó, la transacción
está persistida. En ambos sub-casos: consistente.

En todos los casos, el invariante se mantiene: o la transacción se aplica
completamente, o se puede completar durante la recuperación. \qedhere
\end{proof}

% ====================================================================
\chapter{Arranque y Configuración}
\label{chap:arranque}
% ====================================================================

Este capítulo describe el proceso completo de arranque de \sotos{}, desde el
encendido de la máquina (real o virtual) hasta que el shell LUCAS está listo
para recibir comandos.

% --------------------------------------------------------------------
\section{Bootloader: Limine (BIOS, GPT+FAT32)}
\label{sec:limine}

\sotos{} utiliza el bootloader Limine v8 en modo BIOS. El disco de imagen tiene:
\begin{itemize}
    \item Tabla de particiones GPT con una partición FAT32.
    \item El binario de Limine instalado en el MBR y en los sectores de boot.
    \item La configuración en \texttt{/boot/limine.conf}.
\end{itemize}

La configuración de Limine es mínima:

\begin{lstlisting}[language={},basicstyle=\ttfamily\small,frame=single,caption={Contenido de \texttt{boot/limine.conf}.},label={lst:limine-conf}]
serial: yes
timeout: 0

/sotOS
    protocol: limine
    resolution: 1024x768x32
    kernel_path: boot():/boot/sotos-kernel
    module_path: boot():/boot/initrd.img
\end{lstlisting}

La directiva \texttt{serial: yes} habilita la salida de Limine por el puerto
serial COM1, lo que permite ver los mensajes del bootloader en la terminal QEMU
antes de que el kernel tome control.

% --------------------------------------------------------------------
\section{Secuencia de Arranque del Kernel}
\label{sec:boot-sequence}

La función \texttt{kmain()} ejecuta los siguientes pasos en orden estricto:

\begin{enumerate}
    \item \textbf{Inicializar serial (COM1)}: Primera instrucción. Todas las
          aserciones y mensajes de debug dependen de la salida serial.
          \texttt{kprintln!("sotOS v0.1.0 --- microkernel booting...")}.

    \item \textbf{GDT}: Instalar la Global Descriptor Table con segmentos
          de kernel (Ring~0) y usuario (Ring~3). Cargar el Task State Segment (TSS)
          con pilas de interrupción (IST) para \#DF, \#GP, y misc.

    \item \textbf{SSE}: Habilitar SSE/SSE2 en CR0 y CR4. Crítico para binarios
          musl/glibc que usan instrucciones como \texttt{movaps}. Se aplica al
          BSP aquí y a cada AP durante su inicialización.

    \item \textbf{IDT}: Instalar la Interrupt Descriptor Table con handlers para
          excepciones (\#PF, \#GP, \#DF) e IRQs de hardware.

    \item \textbf{Asignador de marcos}: Inicializar el bitmap desde el mapa de
          memoria del bootloader (vía HHDM).

    \item \textbf{Slab allocator}: Inicializar el heap del kernel (slabs de
          64, 128, 256, 512 bytes).

    \item \textbf{Per-CPU + TSS}: Asignar estructura \texttt{PerCpu} para el BSP,
          recargar GDT con TSS per-CPU (pilas IST únicas por CPU).

    \item \textbf{Capacidades + Planificador}: Inicializar tabla de capacidades,
          pool de objetos, MLQ de 4 niveles. Configurar MSRs de SYSCALL/SYSRET.

    \item \textbf{LAPIC}: Inicializar el Local APIC, calibrar el timer
          ($\sim$100\,Hz periódico), enmascarar el PIT (obsoleto).

    \item \textbf{Proceso init}: Cargar el ELF del initrd, crear primer espacio
          de direcciones de usuario, spawning de los servicios (VMM, kbd, net).
\end{enumerate}

Después de estos pasos, el kernel habilita interrupciones y entra en el bucle
idle (\texttt{hlt} en un loop). Las interrupciones del LAPIC timer causan cambios
de contexto preemptivos al planificador.

% --------------------------------------------------------------------
\section{Bootstrapping del Proceso Init}
\label{sec:init-bootstrap}

El proceso \texttt{init} (cargado desde el initrd CPIO) es el primer código
de espacio de usuario. Su bootstrapping incluye:

\begin{enumerate}
    \item Leer la página \texttt{BootInfo} en \texttt{0xB00000} para obtener
          las capacidades raíz (endpoints, AS caps, frame caps).
    \item Inicializar el framebuffer serial (consola de texto vía COM1).
    \item Registrar el servicio \texttt{"init"} en el registro de servicios del kernel.
    \item Spawning de procesos hijo:
          \begin{itemize}
              \item \textbf{VMM}: Recibe cap de AS del init + cap de notificación.
              \item \textbf{Keyboard}: Recibe cap de IRQ para IRQ~1.
              \item \textbf{Network}: Recibe cap de PCI para escaneo de dispositivos.
          \end{itemize}
    \item Montar el \texttt{ObjectStore} sobre el disco virtio-blk.
    \item Inyectar el sysroot virtual (\texttt{/lib}, \texttt{/bin}, \texttt{/etc},
          \texttt{/tmp}, \texttt{/proc}, \texttt{/sys}).
    \item Forjar el vDSO en \texttt{0xB80000}.
    \item Ejecutar pruebas de boot (hello-linux, musl, busybox, etc.).
    \item Lanzar el shell LUCAS.
\end{enumerate}

% --------------------------------------------------------------------
\section{Layout del Espacio de Direcciones}
\label{sec:addr-layout}

\begin{table}[ht]
\centering
\caption{Mapa completo del espacio de direcciones del proceso init. Todas las
direcciones están en la mitad inferior del espacio virtual (usuario).}
\label{tab:addr-layout}
\begin{tabular}{@{}rll@{}}
\toprule
\textbf{Dirección} & \textbf{Tamaño} & \textbf{Contenido} \\
\midrule
\texttt{0x200000}   & Variable  & ELF \texttt{.text} (init binary) \\
\texttt{0x510000}   & 4\,KiB    & Ring buffer del teclado (compartido) \\
\texttt{0x520000}   & 4\,KiB    & Ring buffer del mouse (compartido) \\
\texttt{0x900000}   & 16\,KiB   & Stack del init (4 páginas, con ASLR) \\
\texttt{0xA00000}   & 4\,KiB    & SPSC ring buffer + stack del producer \\
\texttt{0xB00000}   & 4\,KiB    & BootInfo (solo lectura) \\
\texttt{0xB80000}   & 4\,KiB    & vDSO (R+X, ELF forjado) \\
\texttt{0xC00000+}  & Variable  & Páginas MMIO de virtio (UC) \\
\texttt{0xD00000+}  & 1.25\,MiB & ObjectStore + VFS (320 páginas) \\
\texttt{0xE50000}   & 16\,KiB   & Stack del huésped LUCAS (4 páginas) \\
\texttt{0xE60000}   & 64\,KiB   & Stack del handler LUCAS (16 páginas) \\
\texttt{0xE70000}   & 4\,KiB    & VirtioBlk storage (1 página) \\
\texttt{0xE80000+}  & Variable  & Stacks de procesos hijo (0x2000 c/u) \\
\texttt{0x1000000}  & Variable  & Shell ELF \\
\texttt{0x2000000}  & 1\,MiB    & Heap \texttt{brk} \\
\texttt{0x3000000}  & Variable  & Región \texttt{mmap} \\
\texttt{0x4000000}  & Variable  & Framebuffer (UC, si está presente) \\
\texttt{0x5000000}  & 512\,KiB  & Buffer de spawn (128 páginas) \\
\texttt{0x5200000}  & 128\,KiB  & Buffer de dynamic linking (32 páginas) \\
\texttt{0x5400000}  & 512\,KiB  & Buffer de exec (128 páginas) \\
\texttt{0x6000000}  & Variable  & \texttt{INTERP\_LOAD\_BASE} (ld.so) \\
\texttt{0xA000000}  & Variable  & \texttt{INTERP\_BUF\_BASE} \\
\bottomrule
\end{tabular}
\end{table}

% --------------------------------------------------------------------
\section{Configuración de QEMU}
\label{sec:qemu-config}

El sistema de build utiliza \texttt{just} (make alternativo) con los siguientes
targets principales:

\begin{table}[ht]
\centering
\caption{Targets de \texttt{just} para ejecución en QEMU.}
\label{tab:just-targets}
\begin{tabular}{@{}lp{9cm}@{}}
\toprule
\textbf{Target} & \textbf{Descripción} \\
\midrule
\texttt{just build}    & Compilar el kernel (\texttt{cargo build --package sotos-kernel}) \\
\texttt{just run}      & Build + imagen + QEMU (serial a terminal, sin display, 1 CPU, 1\,GiB RAM) \\
\texttt{just run-gui}  & Como \texttt{run} pero con display SDL + virtio-blk \\
\texttt{just run-net}  & Con virtio-net (SLIRP) + virtio-blk (ObjectStore en disco) \\
\texttt{just run-smp}  & Con 4 CPUs (\texttt{-smp 4}, experimental por race en scheduler) \\
\texttt{just debug}    & Con GDB server (\texttt{-s -S}, conectar con \texttt{target remote :1234}) \\
\bottomrule
\end{tabular}
\end{table}

La línea de QEMU para \texttt{just run-net} (el target más completo) es:

\begin{lstlisting}[language={},basicstyle=\ttfamily\scriptsize,frame=single,breaklines=true,caption={Comando QEMU para \texttt{just run-net}.},label={lst:qemu-cmd}]
qemu-system-x86_64 \
    -drive format=raw,file=target/sotos.img \
    -drive if=none,format=raw,file=target/disk.img,id=disk0 \
    -device virtio-blk-pci,drive=disk0,disable-modern=on \
    -netdev user,id=net0,dns=8.8.8.8,
            hostfwd=udp::5555-:5555,hostfwd=tcp::7777-:7 \
    -device virtio-net-pci,netdev=net0,disable-modern=on \
    -serial stdio -no-reboot -m 1024M
\end{lstlisting}

Aspectos relevantes de la configuración:
\begin{itemize}
    \item \texttt{disable-modern=on}: Fuerza virtio legacy (no MMIO moderno),
          compatible con los drivers de \sotos{}.
    \item \texttt{hostfwd}: Reenvía puertos del host al guest (UDP 5555 para
          pruebas, TCP 7777 para echo server).
    \item \texttt{dns=8.8.8.8}: Configura el DNS del backend SLIRP.
    \item \texttt{-m 1024M}: 1\,GiB de RAM (incrementado desde 256\,MiB original
          para soportar ObjectStore + procesos dinámicos).
    \item \texttt{-serial stdio}: Salida serial del kernel directamente a la terminal.
    \item \texttt{-no-reboot}: QEMU sale en lugar de reiniciar ante triple fault.
\end{itemize}

\section{Resumen de la Parte II}

Esta segunda parte ha cubierto los componentes de espacio de usuario que
transforman las cinco primitivas del micronúcleo en un sistema operativo
funcional:

\begin{enumerate}
    \item \textbf{Gestión de memoria} (\cref{chap:memoria}): Asignador bitmap
          con $\bigO(1)$ amortizado, tablas de páginas de cuatro niveles, y
          Copy-on-Write con pruebas formales de aislamiento y cotas de memoria.

    \item \textbf{LUCAS} (\cref{chap:lucas}): Capa de compatibilidad Linux con
          127+ syscalls emuladas, soporte completo de fork/exec/wait, señales
          asíncronas, y un vDSO forjado para acceso rápido al tiempo.

    \item \textbf{Servicios del sistema} (\cref{chap:servicios}): Red TCP/IP
          (smoltcp, 13 comandos IPC), VMM con CoW, teclado PS/2, y ObjectStore
          transaccional con WAL y snapshots.

    \item \textbf{Arranque} (\cref{chap:arranque}): Limine BIOS, secuencia de
          10 pasos del kernel, bootstrap de init, y configuración de QEMU.
\end{enumerate}

% ====================================================================
% Bibliografía
% ====================================================================
\begin{thebibliography}{99}

\bibitem{blumofe1999}
R.~D. Blumofe and C.~E. Leiserson,
``Scheduling multithreaded computations by work stealing,''
\emph{Journal of the ACM}, vol.~46, no.~5, pp.~720--748, 1999.

\bibitem{liedtke1995}
J.~Liedtke,
``On micro-kernel construction,''
\emph{Proceedings of the 15th ACM Symposium on Operating Systems Principles (SOSP)},
pp.~237--250, 1995.

\bibitem{shapiro1999}
J.~S. Shapiro, J.~M. Smith, and D.~J. Farber,
``EROS: A fast capability system,''
\emph{Proceedings of the 17th ACM Symposium on Operating Systems Principles (SOSP)},
pp.~170--185, 1999.

\bibitem{klein2009}
G.~Klein et~al.,
``seL4: Formal verification of an OS kernel,''
\emph{Proceedings of the 22nd ACM Symposium on Operating Systems Principles (SOSP)},
pp.~207--220, 2009.

\bibitem{engler1995}
D.~R. Engler, M.~F. Kaashoek, and J.~O'Toole,
``Exokernel: An operating system architecture for application-level resource management,''
\emph{Proceedings of the 15th ACM Symposium on Operating Systems Principles (SOSP)},
pp.~251--266, 1995.

\bibitem{dennis1966}
J.~B. Dennis and E.~C. Van~Horn,
``Programming semantics for multiprogrammed computations,''
\emph{Communications of the ACM}, vol.~9, no.~3, pp.~143--155, 1966.

\bibitem{levy1984}
H.~M. Levy,
\emph{Capability-Based Computer Systems},
Digital Press, 1984.

\bibitem{herlihy1991}
M.~Herlihy,
``Wait-free synchronization,''
\emph{ACM Transactions on Programming Languages and Systems}, vol.~13, no.~1,
pp.~124--149, 1991.

\bibitem{cormen2009}
T.~H. Cormen, C.~E. Leiserson, R.~L. Rivest, and C.~Stein,
\emph{Introduction to Algorithms}, 3rd ed.,
MIT Press, 2009.

\end{thebibliography}

\end{document}
